\chapter{Derivatives}
\label{chap:pet-derivatives}

Lie and covariant derivatives act on smooth tensor fields unless a lower
regularity is stated. Unless specified otherwise, $\nabla$ on a Riemannian
manifold denotes its Levi--Civita connection.

\section{Lie Derivatives}

Directional differentiation acts on functions without auxiliary structure;
the flow of a vector field extends it naturally to tensor fields.

\subsection{Directional Derivatives}

\begin{definition}[directional derivative]
  \label{def:pet-ch2-directional-derivative}
  \lean{PetersenLib.directionalDerivative,PetersenLib.gradient}
  \source{petersen:page-0060}
  \leanok
  For $f:M\to\mathbb{R}$ and a
  vector field $Y$, set
  \[
    \nabla_Yf=D_Yf=L_Yf=df(Y)=Y(f).
  \]
  On a Riemannian manifold, the gradient
  $\nabla f=\operatorname{grad}f$ is characterized by
  $g(v,\nabla f)=df(v)$. In coordinates,
  \[
    \nabla f=g^{ij}\partial_i(f)\partial_j,
  \]
  where $(g^{ij})=(g_{ij})^{-1}$.
\end{definition}

\begin{remark}[the gradient is not metric-independent]
  \label{rem:pet-ch2-gradient-not-invariant}
  \lean{PetersenLib.gradient_not_coordinate_invariant}
  \source{petersen:page-0060}
  \uses{def:pet-ch2-directional-derivative}
  \leanok
  The gradient depends on the
  metric. In Cartesian coordinates on $\mathbb{R}^n$,
  $\nabla f=\sum_i\partial_i(f)\partial_i$, but the analogous unweighted
  polar expression
  $\partial_r(f)\partial_r+\partial_\theta(f)\partial_\theta$ is generally
  incorrect. The invariant formulas pair upper and lower indices:
  \[
    df=\partial_i(f)dx^i,\qquad
    g=g_{ij}dx^i dx^j,\qquad
    \nabla f=g^{ij}\partial_i(f)\partial_j.
  \]
\end{remark}

\subsection{Lie Derivatives}

Let $F^t$ denote the local flow of the differentiating vector field.

\begin{definition}[Lie derivative of a function]
  \label{def:pet-ch2-lie-derivative-function}
  \lean{PetersenLib.lieDerivativeFunction}
  \source{petersen:page-0061}
  \leanok
  If $F^t$ is the local flow of
  $X$, define
  \[
    (L_Xf)(p)
      =\lim_{t\to0}\frac{f(F^t(p))-f(p)}t.
  \]
  Equivalently, $f\circ F^t=f+tL_Xf+o(t)$. Thus
  $L_Xf=D_Xf=df(X)$ by
  \cref{def:pet-ch2-directional-derivative}.
\end{definition}

\begin{definition}[Lie derivative of a vector field]
  \label{def:pet-ch2-lie-derivative-vector-field}
  \lean{PetersenLib.lieDerivativeVectorField}
  \source{petersen:page-0061}
  \uses{def:pet-ch2-lie-derivative-function}
  \leanok
  For a vector field $Y$,
  transport $Y_{F^t(p)}$ back to $T_pM$ and define
  \[
    DF^{-t}(Y_{F^t(p)})
      =Y_p+t(L_XY)_p+o(t),
    \qquad
    (L_XY)_p
      =\lim_{t\to0}
         \frac{DF^{-t}(Y_{F^t(p)})-Y_p}{t}.
  \]
  The formalization uses the flow-free manifold bracket
  $\texttt{VectorField.mlieBracket}$. Proposition
  \cref{prop:pet-ch2-lie-derivative-bracket} identifies the two descriptions
  through
  $D_{[X,Y]}f=D_XD_Yf-D_YD_Xf$; the regularity distinction is recorded in
  \cref{rem:pet-ch2-ex-2}.
\end{definition}

\begin{proposition}[Lie derivative of a vector field is the bracket]
  \label{prop:pet-ch2-lie-derivative-bracket}
  \lean{PetersenLib.lieDerivative_vectorField_eq_bracket}
  \source{petersen:page-0061,petersen:page-0062}
  \uses{def:pet-ch2-lie-derivative-vector-field}
  \leanok
  \dcref{ch2:2.1.1}
  If $X,Y$ are vector fields on $M$, then $L_XY=[X,Y]$.
\end{proposition}
\begin{proof}
  \uses{prop:pet-ch2-lie-derivative-bracket}
  \leanok
  \sketch
  The defining expansion gives
  \[
    Y|_{F^t}-DF^t(Y)=t\,DF^t(L_XY)+o(t).
  \]
  Applied to an arbitrary smooth function, its left side equals
  \[
    (D_Yf)\circ F^t-D_Y(f\circ F^t)
      =t(D_XD_Yf-D_YD_Xf)+o(t).
  \]
  Comparison of the first-order terms yields
  $D_{L_XY}f=D_{[X,Y]}f$ for every $f$, hence $L_XY=[X,Y]$.
\end{proof}

\begin{definition}[Lie derivative of a $(0,k)$-tensor]
  \label{def:pet-ch2-lie-derivative-tensor}
  \lean{PetersenLib.lieDerivativeTensor}
  \source{petersen:page-0062}
  \uses{def:pet-ch2-lie-derivative-vector-field}
  \leanok
  For a $(0,k)$-tensor $T$,
  \[
    ((F^t)^*T)(Y_1,\dots,Y_k)
      =T(DF^tY_1,\dots,DF^tY_k)
      =T(Y_1,\dots,Y_k)+t(L_XT)(Y_1,\dots,Y_k)+o(t).
  \]
  This first variation defines $L_XT$.
\end{definition}

\begin{proposition}[Algebraic formula for $L_XT$]
  \label{prop:pet-ch2-lie-derivative-tensor-formula}
  \lean{PetersenLib.lieDerivativeTensor_formula}
  \source{petersen:page-0062,petersen:page-0063}
  \uses{def:pet-ch2-lie-derivative-tensor, prop:pet-ch2-lie-derivative-bracket}
  \leanok
  \dcref{ch2:2.1.2}
  If $X$ is a vector field and $T$ a $(0,k)$-tensor, then
  \[
    (L_XT)(Y_1,\dots,Y_k)
    = D_X\big(T(Y_1,\dots,Y_k)\big)-\sum_{i=1}^kT(Y_1,\dots,L_XY_i,\dots,Y_k).
  \]
\end{proposition}
\begin{proof}
  \uses{prop:pet-ch2-lie-derivative-tensor-formula}
  \leanok
  \sketch
  For $k=1$, substitute
  $Y|_{F^t}=DF^t(Y)+t\,DF^t(L_XY)+o(t)$ in the pullback:
  \[
    ((F^t)^*T)(Y)
      =T(Y)\circ F^t-t\,T(DF^t(L_XY))+o(t).
  \]
  Dividing by $t$ and passing to the limit gives
  $(L_XT)(Y)=D_X(T(Y))-T(L_XY)$. Applying the same expansion in each
  argument proves the stated sum for arbitrary $k$.
\end{proof}


\begin{definition}[product of tensors]
  \label{def:pet-ch2-tensor-product}
  \lean{PetersenLib.TensorOperator.mul}
  \source{petersen:page-0063}
  \uses{def:pet-ch2-lie-derivative-tensor}
  \leanok
  For a $(0,k_1)$-tensor $T_1$ and a $(0,k_2)$-tensor $T_2$, their
  \emph{product} is the $(0,k_1+k_2)$-tensor
  \[
    (T_1\cdot T_2)(X_1,\dots,X_{k_1},Y_1,\dots,Y_{k_2})
    =T_1(X_1,\dots,X_{k_1})\cdot T_2(Y_1,\dots,Y_{k_2}).
  \]
\end{definition}

\begin{proposition}[Product rule]
  \label{prop:pet-ch2-lie-derivative-product-rule}
  \lean{PetersenLib.lieDerivative_product_rule}
  \source{petersen:page-0063}
  \uses{prop:pet-ch2-lie-derivative-tensor-formula}
  \leanok
  \dcref{ch2:2.1.3}
  For $(0,k_i)$-tensors,
  \[
    L_X(T_1\cdot T_2)
      =(L_XT_1)\cdot T_2+T_1\cdot(L_XT_2),
  \]
  where
  \[
    (T_1\cdot T_2)(X_1,\dots,X_{k_1},Y_1,\dots,Y_{k_2})
      =T_1(X_1,\dots,X_{k_1})T_2(Y_1,\dots,Y_{k_2}).
  \]
\end{proposition}
\begin{proof}
  \uses{prop:pet-ch2-lie-derivative-product-rule}
  \leanok
  \sketch
  Evaluate on
  $(X_1,\dots,X_{k_1},Y_1,\dots,Y_{k_2})$. The scalar derivative in
  \cref{prop:pet-ch2-lie-derivative-tensor-formula} satisfies the ordinary
  Leibniz rule, while the slot-correction sum separates into the $X$-slots and
  the $Y$-slots. These two groups are respectively
  $(L_XT_1)\cdot T_2$ and $T_1\cdot(L_XT_2)$.
\end{proof}

\begin{proposition}[Lie derivative in the direction $fX$]
  \label{prop:pet-ch2-lie-derivative-scaling}
  \lean{PetersenLib.lieDerivative_scale_direction}
  \source{petersen:page-0063,petersen:page-0064}
  \uses{prop:pet-ch2-lie-derivative-tensor-formula}
  \leanok
  \dcref{ch2:2.1.4}
  If $T$ is a $(0,k)$-tensor and $f:M\to\mathbb{R}$, then
  \[
    (L_{fX}T)(Y_1,\dots,Y_k)
    = f(L_XT)(Y_1,\dots,Y_k)+\sum_{i=1}^k(L_{Y_i}f)\,T(Y_1,\dots,X,\dots,Y_k),
  \]
  with $X$ replacing $Y_i$ in the $i$th slot.
\end{proposition}
\begin{proof}
  \uses{prop:pet-ch2-lie-derivative-scaling}
  \leanok
  \sketch
  Insert
  $D_{fX}=fD_X$ and
  \[
    [fX,Y_i]=f[X,Y_i]-(D_{Y_i}f)X
  \]
  into \cref{prop:pet-ch2-lie-derivative-tensor-formula}. Multilinearity
  separates the terms multiplied by $f$, which form $fL_XT$, from the
  correction in each slot, giving the asserted sum.
\end{proof}

\begin{lemma}[Locality at a zero of $X$]
  \label{lem:pet-ch2-lie-derivative-local-vanishing}
  \lean{PetersenLib.lieDerivative_local_at_zero}
  \source{petersen:page-0064}
  \uses{prop:pet-ch2-lie-derivative-tensor-formula}
  \leanok
  \dcref{ch2:2.1.5}
  If a vector field $X$ vanishes at $p$, then $(L_XT)|_p$ depends only on the
  value $T|_p$ of the $(0,k)$-tensor $T$ at $p$.
\end{lemma}
\begin{proof}
  \uses{lem:pet-ch2-lie-derivative-local-vanishing}
  \leanok
  \sketch
  At $p$, the first term in
  \[
    (L_XT)(Y_1,\dots,Y_k)
      =D_X(T(Y_1,\dots,Y_k))
       -\sum_iT(Y_1,\dots,L_XY_i,\dots,Y_k)
  \]
  vanishes because $X_p=0$. Every remaining term evaluates $T$ only at $p$,
  so the result depends only on $T_p$.
\end{proof}

\begin{definition}[iterated Lie derivative $(L_XL)_YT$]
  \label{def:pet-ch2-generalized-lie-derivative}
  \lean{PetersenLib.generalizedLieDerivative}
  \source{petersen:page-0064}
  \uses{def:pet-ch2-lie-derivative-tensor}
  \leanok
  For vector fields $X,Y$ and a tensor (or multilinear map on vector fields)
  $T$,
  \[
    (L_XL)_YT := L_X(L_YT)-L_{L_XY}T-L_Y(L_XT) = [L_X,L_Y]T-L_{[X,Y]}T.
  \]
\end{definition}

\begin{proposition}[The Generalized Jacobi Identity]
  \label{prop:pet-ch2-generalized-jacobi-identity}
  \lean{PetersenLib.generalizedJacobiIdentity}
  \source{petersen:page-0065,petersen:page-0066}
  \uses{def:pet-ch2-generalized-lie-derivative}
  \leanok
  \dcref{ch2:2.1.6}
  For all vector fields $X,Y$ and tensors $T$, $(L_XL)_YT=0$.
\end{proposition}
\begin{proof}
  \uses{prop:pet-ch2-generalized-jacobi-identity}
  \leanok
  \sketch
  On functions the identity is
  $[D_X,D_Y]f=D_{[X,Y]}f$. On vector fields it is the Jacobi identity:
  \[
    [X,[Y,Z]]-[Y,[X,Z]]-[[X,Y],Z]=0.
  \]
  For a one-form $\omega$, evaluation on $Z$ gives
  \[
    ([L_X,L_Y]\omega)(Z)
      =[L_X,L_Y](\omega(Z))-\omega([L_X,L_Y]Z).
  \]
  The function and vector-field cases therefore imply
  $([L_X,L_Y]-L_{[X,Y]})\omega=0$. Finally,
  \cref{prop:pet-ch2-lie-derivative-product-rule} extends the result from
  functions, vector fields, and one-forms to their tensor products and sums.
\end{proof}

\begin{definition}[exterior derivative via Lie derivatives]
  \label{def:pet-ch2-exterior-derivative-lie-formula}
  \lean{PetersenLib.exteriorDerivative_lieFormula}
  \source{petersen:page-0066}
  \uses{def:pet-ch2-lie-derivative-tensor}
  \leanok
  For a $k$-form $\omega$,
  \[
    d\omega(X_0,\dots,X_k)
    =\frac12\sum_{i=0}^k(-1)^i(L_{X_i}\omega)(X_0,\dots,\hat X_i,\dots,X_k)
    +\frac12\sum_{i=0}^k(-1)^iL_{X_i}\big(\omega(X_0,\dots,\hat X_i,\dots,X_k)\big).
  \]
  For a $1$-form this specializes to
  $d\omega(X,Y)=D_X(\omega(Y))-D_Y(\omega(X))-\omega([X,Y])$.
\end{definition}

\subsection{Lie Derivatives and the Metric}

\begin{definition}[Hessian via the Lie derivative]
  \label{def:pet-ch2-hessian-lie-derivative}
  \lean{PetersenLib.hessianLieDerivative}
  \source{petersen:page-0066}
  \uses{def:pet-ch2-lie-derivative-tensor, def:pet-ch2-directional-derivative}
  \leanok
  The \emph{Hessian} of $f:(M,g)\to\mathbb{R}$ is the $(0,2)$-tensor
  $\operatorname{Hess}f(X,Y)=\frac12(L_{\nabla f}g)(X,Y)$.
\end{definition}

\begin{proposition}[Hessian at a critical point is metric-independent]
  \label{prop:pet-ch2-hessian-critical-point}
  \lean{PetersenLib.hessian_criticalPoint_metricIndependent}
  \source{petersen:page-0066,petersen:page-0067}
  \uses{def:pet-ch2-hessian-lie-derivative}
  \leanok
  If $df|_p=0$, choose coordinates $x^i$ around $p$ with $g_{ij}|_p=\delta_{ij}$.
  Then $\operatorname{Hess}f|_p$ computed from $g$ equals $\operatorname{Hess}f|_p$
  computed from the flat metric $\delta_{ij}dx^idx^j$ in the same coordinates.
\end{proposition}
\begin{proof}
  \leanok
  \uses{lem:pet-ch2-lie-derivative-local-vanishing}
  \sketch
  Since $df_p=0$, the vector field
  $\nabla f$ vanishes at $p$. By
  \cref{lem:pet-ch2-lie-derivative-local-vanishing},
  $(L_{\nabla f}g)_p$ depends only on $g_p$. In coordinates with
  $g_{ij}(p)=\delta_{ij}$ it therefore equals
  $L_{\nabla f}(\delta_{ij}dx^i dx^j)|_p$, proving the claim after division
  by $2$.
\end{proof}

\begin{definition}[divergence via the Lie derivative]
  \label{def:pet-ch2-divergence-lie-derivative}
  \lean{PetersenLib.divergenceLieDerivative}
  \source{petersen:page-0067}
  \uses{def:pet-ch2-lie-derivative-tensor}
  \leanok
  The \emph{divergence} of a vector field $X$ is the function $\operatorname{div}X$
  measuring how the volume form changes along the flow of $X$:
  $L_X\,\mathrm{vol}=(\operatorname{div}X)\,\mathrm{vol}$.
\end{definition}

\begin{remark}[$L_X\mathrm{vol}$ is exact]
  \label{rem:pet-ch2-lie-derivative-volume-exact}
  \lean{PetersenLib.lieDerivative_volume_exact}
  \source{petersen:page-0067}
  \uses{def:pet-ch2-divergence-lie-derivative}
  \leanok
  Interior multiplication in
  the first variable satisfies
  \[
    L_X\mathrm{vol}=d(i_X\mathrm{vol}).
  \]
  \sketch
  Cartan's identity $L_X=d\,i_X+i_Xd$ reduces to the formula because
  $d\,\mathrm{vol}=0$ in top degree. Evaluating on a positively oriented
  orthonormal frame gives
  $L_X\mathrm{vol}=(\operatorname{div}X)\mathrm{vol}$.
\end{remark}

\begin{definition}[Laplacian]
  \label{def:pet-ch2-laplacian}
  \lean{PetersenLib.laplacian}
  \source{petersen:page-0067}
  \uses{def:pet-ch2-divergence-lie-derivative, def:pet-ch2-directional-derivative}
  \leanok
  The \emph{Laplacian} of $f$ is $\Delta f=\operatorname{div}\nabla f$.
\end{definition}

\begin{proposition}[divergence as a trace of $L_Xg$]
  \label{prop:pet-ch2-divergence-trace-formula}
  \lean{PetersenLib.divergence_trace_formula}
  \source{petersen:page-0067}
  \uses{def:pet-ch2-divergence-lie-derivative}
  \leanok
  For a positively oriented orthonormal frame $E_1,\dots,E_n$,
  $\operatorname{div}X=\sum_i\frac12(L_Xg)(E_i,E_i)$; consequently
  $\Delta f=\operatorname{div}\nabla f=\sum_i\operatorname{Hess}f(E_i,E_i)$.
\end{proposition}
\begin{proof}
  \uses{prop:pet-ch2-divergence-trace-formula}
  \leanok
  \sketch
  Choose an oriented orthonormal
  frame $E_1,\dots,E_n$. Since
  $\mathrm{vol}(E_1,\dots,E_n)=1$,
  \[
    \operatorname{div}X
      =(L_X\mathrm{vol})(E_1,\dots,E_n)
      =-\sum_i g(L_XE_i,E_i).
  \]
  Because $g(E_i,E_i)=1$,
  $(L_Xg)(E_i,E_i)=-2g(L_XE_i,E_i)$. Summation gives the divergence
  formula; applying it to $X=\nabla f$ gives the Laplacian formula.
\end{proof}

\begin{proposition}[Hessian agrees with the Euclidean Hessian]
  \label{prop:pet-ch2-hessian-euclidean-agreement}
  \lean{PetersenLib.hessian_euclidean_agreement}
  \source{petersen:page-0067,petersen:page-0068}
  \uses{def:pet-ch2-hessian-lie-derivative}
  \leanok
  For $f:\mathbb{R}^n\to\mathbb{R}$ with the canonical metric, the Hessian of
  \cref{def:pet-ch2-hessian-lie-derivative} equals the classical Hessian:
  $L_{\nabla f}(\delta_{ij}dx^idx^j)=2\,\partial_{ij}f\,dx^idx^j$, i.e.
  $\operatorname{Hess}f=\partial_{ij}f\,dx^idx^j$.
\end{proposition}
\begin{proof}
  \uses{prop:pet-ch2-hessian-euclidean-agreement}
  \leanok
  \sketch
  In Cartesian coordinates,
  $\nabla f=(\partial_i f)\partial_i$ and
  $L_{\nabla f}dx^i=d(\partial_i f)$. Hence
  \[
    L_{\nabla f}g
      =\delta_{ij}\big(d(\partial_i f)\,dx^j
                       +dx^i\,d(\partial_j f)\big)
      =2\,\partial_{ij}f\,dx^i dx^j.
  \]
  The conclusion follows from
  $\operatorname{Hess}f=\tfrac12L_{\nabla f}g$.
\end{proof}

\subsection{Lie Groups}

\begin{lemma}[Differential of $\operatorname{Ad}$]
  \label{lem:pet-ch2-ad-differential}
  \lean{PetersenLib.mfderiv_adMap_apply_eq_groupBracket}
  \source{petersen:page-0068}
  \leanok
  \dcref{ch2:2.1.7}
  For a Lie group $G$ with inner automorphism $\operatorname{Ad}_h:x\mapsto hxh^{-1}$
  and its differential $\operatorname{Ad}_h:\mathfrak g\to\mathfrak g$ at $e$, the
  differential of $h\mapsto\operatorname{Ad}_h$ is $U\mapsto\operatorname{ad}_U(X)=[U,X]$.
\end{lemma}
\begin{proof}
  \uses{lem:pet-ch2-ad-differential}
  \leanok
  \sketch
  For the flow $F^t$ of a left-invariant
  field $U$, left invariance gives $F^t(x)=xF^t(e)$ and
  $DF^t=DR_{F^t(e)}$. Therefore
  \[
    \frac d{dt}\operatorname{Ad}_{F^t(e)}(X)\Big|_{t=0}
      =\frac d{dt}\big(DF^{-t}(X|_{F^t(e)})\big)\Big|_{t=0}
      =L_UX=[U,X],
  \]
  by \cref{def:pet-ch2-lie-derivative-vector-field} and
  \cref{prop:pet-ch2-lie-derivative-bracket}.
\end{proof}

\begin{lemma}[Bracket on $\mathfrak{gl}(V)$]
  \label{lem:pet-ch2-gl-bracket-commutator}
  \lean{PetersenLib.gl_bracket_eq_commutator}
  \source{petersen:page-0068,petersen:page-0069}
  \uses{lem:pet-ch2-ad-differential}
  \leanok
  \dcref{ch2:2.1.8}
  Let $G=\mathrm{GL}(V)$. The Lie bracket on $\mathfrak{gl}(V)=T_I\mathrm{GL}(V)$
  of left-invariant vector fields is given by commutation of linear maps:
  $[X,Y]_{\mathfrak l}=XY-YX$.
\end{lemma}
\begin{proof}
  \leanok
  \uses{lem:pet-ch2-gl-bracket-commutator}
  \sketch
  Here
  $\operatorname{Ad}_h(X)=hXh^{-1}$. With
  $F^t(I)=I+tU+o(t)$, \cref{lem:pet-ch2-ad-differential} gives
  \[
    [U,X]
      =\frac d{dt}\big((I+tU)X(I-tU)\big)\Big|_{t=0}
      =UX-XU.
  \]
\end{proof}

\section{Connections}

Connections replace coordinatewise differentiation by an operator satisfying
tensoriality in the direction variable and the Leibniz rule in the field.

\subsection{Covariant Differentiation}

\begin{definition}[covariant derivative on $\mathbb{R}^n$]
  \label{def:pet-ch2-covariant-derivative-euclidean}
  \lean{PetersenLib.covariantDerivativeEuclidean}
  \source{petersen:page-0069}
  \uses{def:pet-ch2-directional-derivative}
  \leanok
  In Cartesian
  coordinates on $\mathbb{R}^n$, for $X=X^i\partial_i$ define
  \[
    \nabla_YX=dX^i(Y)\partial_i.
  \]
  Thus constant-coefficient fields are parallel. The formula is tied to the
  Cartesian frame: for example
  $\partial_r=r^{-1}(x\partial_x+y\partial_y)$ and
  $\partial_\theta=-y\partial_x+x\partial_y$ have nonconstant Cartesian
  coefficients.
\end{definition}

\begin{definition}[dual $1$-form $\theta_X$]
  \label{def:pet-ch2-dual-one-form}
  \lean{PetersenLib.dualOneForm}
  \source{petersen:page-0069}
  \leanok
  For a vector field $X$ on $(M,g)$, the dual $1$-form is
  $\theta_X:=i_Xg$, i.e.\ $\theta_X(Y)=g(X,Y)$. $X$ is locally a gradient field
  iff $d\theta_X=0$; in general $d\theta_X$ measures the extent to which $X$
  fails to be a gradient field.
\end{definition}

\begin{definition}[Killing field]
  \label{def:pet-ch2-killing-field}
  \lean{PetersenLib.IsKillingField}
  \source{petersen:page-0069}
  \uses{def:pet-ch2-lie-derivative-tensor}
  \leanok
  A vector field $X$ on $(M,g)$ is a \emph{Killing field} if $L_Xg=0$;
  equivalently, if $F^t$ is the local flow of $X$, the $F^t$ are local
  isometries.
\end{definition}

\begin{remark}[Killing fields on $\mathbb{R}^n$]
  \label{rem:pet-ch2-killing-field-euclidean}
  \lean{PetersenLib.killingField_euclidean_characterization}
  \source{petersen:page-0070}
  \uses{def:pet-ch2-killing-field}
  \leanok
  A field
  $X=X^i\partial_i$ on $\mathbb{R}^n$ is Killing exactly when
  \[
    \partial_iX^j+\partial_jX^i=0.
  \]
  \sketch
  Differentiating and cyclically permuting the indices forces all second
  derivatives of the components to vanish. Hence
  \[
    X(x)=Ax+\beta,\qquad A^{\mathsf T}=-A.
  \]
  Its stated flow is $F^t(x)=\exp(At)x+t\beta$; rotations include
  $\partial_\theta$ on $\mathbb{R}^2$. A Euclidean field is constant exactly
  when it is both Killing and locally a gradient.
\end{remark}

\begin{proposition}[Implicit formula on $\mathbb{R}^n$]
  \label{prop:pet-ch2-covariant-derivative-implicit-euclidean}
  \lean{PetersenLib.covariantDerivative_implicit_euclidean}
  \source{petersen:page-0070,petersen:page-0071}
  \uses{def:pet-ch2-covariant-derivative-euclidean, def:pet-ch2-dual-one-form}
  \leanok
  \dcref{ch2:2.2.1}
  The covariant derivative on $\mathbb{R}^n$ satisfies
  $2g(\nabla_YX,Z)=(L_Xg)(Y,Z)+(d\theta_X)(Y,Z)$.
\end{proposition}
\begin{proof}
  \uses{prop:pet-ch2-covariant-derivative-implicit-euclidean}
  \leanok
  \sketch
  Both sides are
  tensorial in $Y,Z$, so take $Y=\partial_k$, $Z=\partial_\ell$ and write
  $X=a^i\partial_i$. Since
  $L_X\partial_k=-(\partial_k a^i)\partial_i$,
  \[
    (L_Xg)(\partial_k,\partial_\ell)
       =\partial_k a^\ell+\partial_\ell a^k,\qquad
    d\theta_X(\partial_k,\partial_\ell)
       =\partial_k a^\ell-\partial_\ell a^k.
  \]
  Their sum is $2\partial_k a^\ell
  =2g(\nabla_{\partial_k}X,\partial_\ell)$.
\end{proof}

\begin{definition}[affine connection]
  \label{def:pet-ch2-affine-connection}
  \lean{PetersenLib.AffineConnection}
  \source{petersen:page-0071,petersen:page-0073}
  \leanok
  An \emph{affine connection} assigns
  $\nabla_vX\in T_pM$ to $v\in T_pM$ and a vector field $X$, subject to
  \[
    \nabla_{\alpha v+\beta w}X
      =\alpha\nabla_vX+\beta\nabla_wX,
    \qquad
    \nabla_Y(fX)=(D_Yf)X+f\nabla_YX,
  \]
  and additivity in $X$. Thus the direction variable is tensorial and the
  differentiated-field variable satisfies the Leibniz rule.
\end{definition}

\begin{definition}[Riemannian connection]
  \label{def:pet-ch2-riemannian-connection}
  \lean{PetersenLib.RiemannianConnection}
  \source{petersen:page-0071,petersen:page-0073}
  \uses{def:pet-ch2-affine-connection}
  \leanok
  An affine connection in the sense of
  \cref{def:pet-ch2-affine-connection} on $(M,g)$
  is \emph{Riemannian} if it is torsion free and metric-compatible:
  \[
    \nabla_XY-\nabla_YX=[X,Y],
    \qquad
    D_Zg(X,Y)=g(\nabla_ZX,Y)+g(X,\nabla_ZY).
  \]
\end{definition}

\begin{remark}[Koszul's formula]
  \label{rem:pet-ch2-koszul-formula}
  \lean{PetersenLib.koszulFormula}
  \source{petersen:page-0072}
  \uses{prop:pet-ch2-covariant-derivative-implicit-euclidean, def:pet-ch2-dual-one-form, def:pet-ch2-lie-derivative-tensor}
  \leanok
  Applying the tensor Lie-derivative formula
  \cref{prop:pet-ch2-lie-derivative-tensor-formula} and the exterior-derivative
  formula \cref{def:pet-ch2-exterior-derivative-lie-formula} to the Euclidean
  prototype \cref{prop:pet-ch2-covariant-derivative-implicit-euclidean} gives
  \[
  \begin{aligned}
    (L_Xg)(Y,Z)+(d\theta_X)(Y,Z)
      ={}&D_Xg(Y,Z)+D_Yg(X,Z)-D_Zg(X,Y)\\
         &-g([X,Y],Z)-g([Y,Z],X)+g([Z,X],Y).
  \end{aligned}
  \]
  This is the Koszul expression determining the Riemannian connection.
\end{remark}


\begin{lemma}[extension of a tangent vector to a vector field]
  \label{lem:pet-ch2-extend-tangent-vector}
  \lean{PetersenLib.exists_smoothVectorField_eq}
  \leanok
  On a smooth ($\sigma$-compact, Hausdorff) manifold, every tangent vector
  $v\in T_pM$ is the value at $p$ of a globally defined smooth vector field.
\end{lemma}
\begin{proof}
  \leanok
  \sketch
  Extend $v$ as a constant-coefficient
  field in a chart at $p$, then multiply by a smooth bump function supported
  in the chart and equal to $1$ near $p$. Extension by zero is a global
  smooth field with value $v$ at $p$.
\end{proof}

\begin{lemma}[Riemannian connections satisfy Koszul's formula]
  \label{lem:pet-ch2-riemannian-connection-koszul}
  \lean{PetersenLib.RiemannianConnection.koszul}
  \uses{def:pet-ch2-riemannian-connection, rem:pet-ch2-koszul-formula}
  \leanok
  Every Riemannian connection $\nabla$ on $(M,g)$ satisfies
  $2g(\nabla_YX,Z)=D_Xg(Y,Z)+D_Yg(X,Z)-D_Zg(X,Y)
  -g([X,Y],Z)-g([Y,Z],X)+g([Z,X],Y)$.
\end{lemma}
\begin{proof}
  \leanok
  \sketch
  Expand the three metric
  derivatives by compatibility and replace every bracket by the
  torsion-free identity. Symmetry of $g$ cancels all terms except
  $2g(\nabla_YX,Z)$.
\end{proof}

\begin{theorem}[The Fundamental Theorem of Riemannian Geometry]
  \label{thm:pet-ch2-fundamental-theorem}
  \lean{PetersenLib.fundamentalTheoremRiemannianGeometry}
  \source{petersen:page-0071,petersen:page-0072,petersen:page-0073}
  \uses{def:pet-ch2-riemannian-connection, rem:pet-ch2-koszul-formula, lem:pet-ch2-extend-tangent-vector, lem:pet-ch2-riemannian-connection-koszul}
  \leanok
  \dcref{ch2:2.2.2}
  On $(M,g)$ there is exactly one Riemannian connection $\nabla$
  (\cref{def:pet-ch2-riemannian-connection}), given implicitly by Koszul's
  formula (\cref{rem:pet-ch2-koszul-formula}):
  \[
    2g(\nabla_YX,Z)=D_Xg(Y,Z)+D_Yg(X,Z)-D_Zg(X,Y)
    -g([X,Y],Z)-g([Y,Z],X)+g([Z,X],Y).
  \]
\end{theorem}
\begin{proof}
  \uses{thm:pet-ch2-fundamental-theorem}
  \leanok
  \sketch
  Define $\nabla_YX$ by the Koszul
  expression. Nondegeneracy of $g$ determines it uniquely. The expression is
  linear and tensorial in $Y$, and
  \[
    L_{fX}g+d\theta_{fX}
      =f(L_Xg+d\theta_X)+2\,df\cdot\theta_X
  \]
  gives $\nabla_Y(fX)=f\nabla_YX+(D_Yf)X$. Antisymmetrizing the Koszul
  expression in $X,Y$ yields
  $\nabla_XY-\nabla_YX=[X,Y]$; symmetrizing the corresponding formulas yields
  \[
    D_Zg(X,Y)=g(\nabla_ZX,Y)+g(X,\nabla_ZY).
  \]
  Thus the constructed connection is Riemannian. Conversely, expanding the
  right side of Koszul's formula using metric compatibility and
  torsion-freeness reduces it to $2g(\nabla_YX,Z)$, so any Riemannian
  connection agrees with the construction.
\end{proof}

\begin{lemma}[Locality on open sets]
  \label{lem:pet-ch2-locality-open-set}
  \lean{PetersenLib.connection_local_openSet}
  \source{petersen:page-0073,petersen:page-0074}
  \uses{def:pet-ch2-affine-connection}
  \leanok
  \dcref{ch2:2.2.3}
  Let $\nabla$ be an affine connection on $M$, $p\in M$, $v\in T_pM$, and
  $X,Y$ vector fields with $X=Y$ on a neighborhood $U\ni p$. Then
  $\nabla_vX=\nabla_vY$.
\end{lemma}
\begin{proof}
  \uses{lem:pet-ch2-locality-open-set}
  \leanok
  \sketch
  Choose a smooth $\lambda$ supported in
  $U$ and equal to $1$ near $p$. Then $\lambda X=\lambda Y$ globally, while
  $\lambda(p)=1$ and $d\lambda_p=0$ imply
  $\nabla_v(\lambda X)=\nabla_vX$ and
  $\nabla_v(\lambda Y)=\nabla_vY$.
\end{proof}

\begin{lemma}[Locality along a curve]
  \label{lem:pet-ch2-locality-along-curve}
  \lean{PetersenLib.connection_local_alongCurve}
  \source{petersen:page-0074}
  \uses{def:pet-ch2-affine-connection}
  \leanok
  \dcref{ch2:2.2.4}
  Let $\nabla$ be an affine connection, $X$ a vector field, and $c:I\to M$
  smooth with $c(0)=p$ and $\dot c(0)=v\in T_pM$. If $X\circ c=Y\circ c$
  for another vector field $Y$, then $\nabla_vX=\nabla_vY$.
\end{lemma}
\begin{proof}
  \leanok
  \uses{def:pet-ch2-affine-connection}
  \sketch
  In a local frame $E_i$, write
  $X=X^iE_i$ and $Y=Y^iE_i$. Equality along $c$ gives
  $X^i(p)=Y^i(p)$ and
  $dX^i(v)=\frac d{dt}(X^i\circ c)|_0
          =\frac d{dt}(Y^i\circ c)|_0=dY^i(v)$.
  The connection's Leibniz rule applied to the two frame expansions then
  gives $\nabla_vX=\nabla_vY$.
\end{proof}

\begin{definition}[normal coordinates and normal frame at $p$]
  \label{def:pet-ch2-normal-coordinates-frame}
  \lean{PetersenLib.isNormalCoordinatesAt,PetersenLib.isNormalFrameAt}
  \source{petersen:page-0074}
  \uses{def:pet-ch2-riemannian-connection}
  \leanok
  A coordinate system is \emph{normal} at $p$ if $g_{ij}|_p=\delta_{ij}$ and
  $\partial_kg_{ij}|_p=0$. An orthonormal frame $E_1,\dots,E_n$ is \emph{normal}
  at $p$ if $(\nabla_vE_i)|_p=0$ for all $i$ and $v\in T_pM$. Such coordinates
  and frames always exist at every point.
\end{definition}

\subsection{Covariant Derivatives of Tensors}

The Levi--Civita connection extends uniquely to the tensor algebra by the
Leibniz rule and compatibility with contractions.

\begin{definition}[covariant derivative of an $(s,t)$-tensor field]
  \label{def:pet-ch2-covariant-derivative-tensor}
  \lean{PetersenLib.covariantDerivativeTensor}
  \source{petersen:page-0075}
  \uses{def:pet-ch2-riemannian-connection}
  \leanok
  For an $(s,t)$-tensor field $S$, the \emph{covariant derivative} $\nabla S$
  is the $(s,t+1)$-tensor field determined by requiring the Leibniz rule
  $\nabla_X(S(Y))=(\nabla_XS)(Y)+S(\nabla_XY)$: for $s=0,1$,
  \[
    \nabla S(X,Y_1,\dots,Y_r)=(\nabla_XS)(Y_1,\dots,Y_r)
    =\nabla_X\big(S(Y_1,\dots,Y_r)\big)-\sum_{i=1}^rS(Y_1,\dots,\nabla_XY_i,\dots,Y_r),
  \]
  where $\nabla_X$ acts as $D_X$ on functions and as covariant differentiation
  on vector fields; equivalently $\nabla_XS=[\nabla_X,S]$ when $S$ is
  $(1,1)$. For general $s$, extend by the product rule
  $\nabla_X(X_1\otimes X_2)=(\nabla_XX_1)\otimes X_2+X_1\otimes(\nabla_XX_2)$.
\end{definition}

\begin{definition}[parallel tensor]
  \label{def:pet-ch2-parallel-tensor}
  \lean{PetersenLib.IsParallel}
  \source{petersen:page-0075}
  \uses{def:pet-ch2-covariant-derivative-tensor}
  \leanok
  A tensor $S$ is \emph{parallel} if $\nabla S=0$. On Euclidean space, a
  tensor written in Cartesian coordinates is parallel iff it has constant
  coefficients.
\end{definition}

\begin{proposition}[The metric and volume form are parallel]
  \label{prop:pet-ch2-metric-volume-parallel}
  \lean{PetersenLib.metric_and_volume_parallel}
  \source{petersen:page-0076}
  \uses{def:pet-ch2-parallel-tensor}
  \leanok
  \dcref{ch2:2.2.5}
  On a Riemannian $n$-manifold, $\nabla g=0$ and $\nabla\mathrm{vol}=0$.
\end{proposition}
\begin{proof}
  \uses{prop:pet-ch2-metric-volume-parallel}
  \leanok
  \sketch
  Metric compatibility is precisely
  $\nabla g=0$. For an oriented orthonormal frame,
  \[
    (\nabla_X\mathrm{vol})(E_1,\dots,E_n)
      =-\sum_i\mathrm{vol}(E_1,\dots,\nabla_XE_i,\dots,E_n)
      =-\sum_i g(\nabla_XE_i,E_i)=0,
  \]
  because each $g(E_i,E_i)$ is constant.
\end{proof}

\begin{proposition}[Hessian via the covariant derivative]
  \label{prop:pet-ch2-hessian-covariant-derivative}
  \lean{PetersenLib.hessian_via_covariantDerivative}
  \source{petersen:page-0076,petersen:page-0077}
  \uses{def:pet-ch2-covariant-derivative-tensor, def:pet-ch2-hessian-lie-derivative}
  \leanok
  \dcref{ch2:2.2.6}
  For $f:(M,g)\to\mathbb{R}$, $(\nabla df)(X,Y)=g(\nabla_X\nabla f,Y)=\operatorname{Hess}f(X,Y)$.
\end{proposition}
\begin{proof}
  \uses{prop:pet-ch2-hessian-covariant-derivative}
  \leanok
  \sketch
  Directly from the tensor derivative of
  \cref{def:pet-ch2-covariant-derivative-tensor},
  \[
    (\nabla_Xdf)(Y)=D_XD_Yf-D_{\nabla_XY}f.
  \]
  Torsion-freeness makes this expression symmetric in $X,Y$. Metric
  compatibility also gives
  $(\nabla_Xdf)(Y)=g(\nabla_X\nabla f,Y)$. Expanding
  $\tfrac12L_{\nabla f}g$ and replacing brackets by torsion-free covariant
  derivatives produces the same symmetric expression, hence all three
  tensors agree.
\end{proof}

\subsection{The Adjoint of the Covariant Derivative}

\begin{definition}[adjoint of the covariant derivative]
  \label{def:pet-ch2-covariant-derivative-adjoint}
  \lean{PetersenLib.covariantDerivativeAdjoint}
  \source{petersen:page-0077}
  \uses{def:pet-ch2-covariant-derivative-tensor}
  \leanok
  For an $(s,t)$-tensor $S$ with $t>0$, and an orthonormal frame
  $E_1,\dots,E_n$, the \emph{adjoint} $\nabla^*S$ is the $(s,t-1)$-tensor
  \[
    (\nabla^*S)(X_2,\dots,X_t)=-\sum_i(\nabla_{E_i}S)(E_i,X_2,\dots,X_t).
  \]
\end{definition}

\begin{proposition}[Divergence as adjoint of $\theta_X$]
  \label{prop:pet-ch2-divergence-adjoint-formula}
  \lean{PetersenLib.divergence_eq_neg_adjoint_dualForm}
  \source{petersen:page-0077}
  \uses{def:pet-ch2-covariant-derivative-adjoint, def:pet-ch2-divergence-lie-derivative, def:pet-ch2-dual-one-form}
  \leanok
  \dcref{ch2:2.2.7}
  If $X$ is a vector field with dual $1$-form $\theta_X$, then
  $\operatorname{div}X=-\nabla^*\theta_X$.
\end{proposition}
\begin{proof}
  \uses{prop:pet-ch2-divergence-adjoint-formula}
  \leanok
  \sketch
  In an orthonormal frame,
  \[
    -\nabla^*\theta_X
      =\sum_i(\nabla_{E_i}\theta_X)(E_i)
      =\sum_i g(\nabla_{E_i}X,E_i)
      =\sum_i\tfrac12(L_Xg)(E_i,E_i).
  \]
  The final sum is $\operatorname{div}X$ by
  \cref{prop:pet-ch2-divergence-trace-formula}.
\end{proof}

\begin{proposition}[$\nabla^*$ is the $L^2$-adjoint of $\nabla$]
  \label{prop:pet-ch2-adjoint-integral-identity}
  \lean{PetersenLib.adjoint_L2_identity}
  \source{petersen:page-0078}
  \uses{def:pet-ch2-covariant-derivative-adjoint, def:pet-ch2-covariant-derivative-tensor}
  \leanok
  \dcref{ch2:2.2.8}
  If $S$ is a compactly supported $(s,t)$-tensor and $T$ a compactly supported
  $(s,t+1)$-tensor, then $\int g(\nabla S,T)\,\mathrm{vol}=\int g(S,\nabla^*T)\,\mathrm{vol}$.
\end{proposition}
\begin{proof}
  \uses{prop:pet-ch2-adjoint-integral-identity}
  \sketch
  For $s=t=1$, set
  $\omega(X)=g(i_XT,S)$. In an orthonormal frame, the tensor-product
  Leibniz rule and metric compatibility give the pointwise Green identity
  \[
    -\nabla^*\omega=g(\nabla S,T)-g(S,\nabla^*T).
  \]
  The connection terms from differentiating the frame cancel by
  $g(\nabla_XE_a,E_b)=-g(E_a,\nabla_XE_b)$; equivalently one may evaluate in a
  normal frame at each point. Integrating and applying the compact-support
  divergence theorem to $\omega^\sharp$, using
  \cref{rem:pet-ch2-lie-derivative-volume-exact}, proves the formula. The
  formalized integral statement takes this divergence theorem as an explicit
  hypothesis, since the Riemannian volume measure discussed in
  \cref{rem:pet-ch1-integration-riemannian} is not yet available; the pointwise
  identity is unconditional and extends to all tensor types by the same
  contraction calculation.
\end{proof}

\subsection{Exterior Derivatives}

\begin{definition}[exterior derivative via the covariant derivative]
  \label{def:pet-ch2-exterior-derivative-covariant}
  \lean{PetersenLib.exteriorDerivative_covariantFormula}
  \source{petersen:page-0078,petersen:page-0079}
  \uses{def:pet-ch2-covariant-derivative-tensor}
  \leanok
  For a $k$-form $\omega$, $(d\omega)(X_0,\dots,X_k)=\sum_i(-1)^i(\nabla_{X_i}\omega)(X_0,\dots,\hat X_i,\dots,X_k)$.
  More generally, for a $(1,k)$-tensor $T$ skew-symmetric in its $k$
  covariant slots, define the $(1,k+1)$-tensor
  $(d^\nabla T)(X_0,\dots,X_k)=\sum_i(-1)^i(\nabla_{X_i}T)(X_0,\dots,\hat X_i,\dots,X_k)$;
  for $k=0$, $T=Y$ a vector field, $(d^\nabla Y)(X)=\nabla_XY$; for $k=1$,
  \[
    (d^\nabla T)(X,Y)=(\nabla_XT)(Y)-(\nabla_YT)(X)
    =\nabla_X(T(Y))-\nabla_Y(T(X))-T([X,Y]).
  \]
\end{definition}

\subsection{The Second Covariant Derivative}

\begin{definition}[second covariant derivative]
  \label{def:pet-ch2-second-covariant-derivative}
  \lean{PetersenLib.secondCovariantDerivative}
  \source{petersen:page-0079}
  \uses{def:pet-ch2-covariant-derivative-tensor}
  \leanok
  For an $(s,t)$-tensor field $S$, the \emph{second covariant derivative}
  $\nabla^2S$ is the $(s,t+2)$-tensor field
  \[
    (\nabla^2_{X_1,X_2}S)(Y_1,\dots,Y_t)
    = (\nabla_{X_1}(\nabla S))(X_2,Y_1,\dots,Y_t)
    = (\nabla_{X_1}(\nabla_{X_2}S))(Y_1,\dots,Y_t)-(\nabla_{\nabla_{X_1}X_2}S)(Y_1,\dots,Y_t).
  \]
\end{definition}

\begin{remark}[Hessian as second covariant derivative; Laplacian as trace]
  \label{rem:pet-ch2-hessian-second-covariant-derivative}
  \lean{PetersenLib.hessian_eq_secondCovariantDerivative}
  \source{petersen:page-0079}
  \uses{def:pet-ch2-second-covariant-derivative, def:pet-ch2-laplacian, def:pet-ch2-covariant-derivative-adjoint}
  \leanok
  For $f:M\to\mathbb{R}$, $\nabla^2_{X,Y}f=\nabla_X\nabla_Yf-\nabla_{\nabla_XY}f
  =\nabla_Xdf(Y)-df(\nabla_XY)=(\nabla_Xdf)(Y)=\operatorname{Hess}f(X,Y)$, and
  $\nabla^2f$ is symmetric in $X,Y$ (unlike the second covariant derivative of
  a general tensor, whose failure of symmetry is central to curvature).
  Consequently $\Delta f=-\nabla^*\nabla f=\sum_i\nabla^2_{E_i,E_i}f$.
\end{remark}

\subsection{The Lie Derivative of the Covariant Derivative}

\begin{definition}[Lie derivative of the connection]
  \label{def:pet-ch2-lie-derivative-connection}
  \lean{PetersenLib.lieDerivativeConnection}
  \source{petersen:page-0079,petersen:page-0080}
  \uses{def:pet-ch2-lie-derivative-vector-field, def:pet-ch2-riemannian-connection}
  \leanok
  \[
    (L_X\nabla)_UV := (L_X\nabla)(U,V)
    = L_X(\nabla_UV)-\nabla_{L_XU}V-\nabla_UL_XV
    = [X,\nabla_UV]-\nabla_{[X,U]}V-\nabla_U[X,V].
  \]
\end{definition}

\begin{remark}[symmetry and tensoriality of $L_X\nabla$]
  \label{rem:pet-ch2-lie-derivative-connection-symmetric}
  \lean{PetersenLib.lieDerivativeConnection_symmetric}
  \source{petersen:page-0080}
  \uses{def:pet-ch2-lie-derivative-connection}
  \leanok
  $(L_X\nabla)(U,V)$ is tensorial in $U$ (since $\nabla_UV$ is tensorial in
  $U$), and it is symmetric in $U,V$; hence it is tensorial in both
  variables.
\end{remark}
\begin{proof}
  \uses{rem:pet-ch2-lie-derivative-connection-symmetric}
  \leanok
  \sketch
  Tensoriality in $U$
  follows directly from the definition. Torsion-freeness and
  \cref{prop:pet-ch2-lie-derivative-bracket} identify the antisymmetric part as
  \[
    (L_X\nabla)(U,V)-(L_X\nabla)(V,U)
      =L_X(L_UV)-L_{L_XU}V-L_U(L_XV)
      =(L_XL)_UV.
  \]
  This vanishes by
  \cref{prop:pet-ch2-generalized-jacobi-identity}. Symmetry then transfers
  tensoriality from the first argument to the second.
\end{proof}

\subsection{The Covariant Derivative of the Covariant Derivative}

\begin{definition}[covariant derivative of the connection]
  \label{def:pet-ch2-covariant-derivative-connection}
  \lean{PetersenLib.covariantDerivativeConnection}
  \source{petersen:page-0080}
  \uses{def:pet-ch2-covariant-derivative-tensor, def:pet-ch2-second-covariant-derivative}
  \leanok
  \[
    (\nabla_X\nabla)_YT := \nabla_X(\nabla_YT)-\nabla_{\nabla_XY}T-\nabla_Y(\nabla_XT).
  \]
  This is not tensorial in $X$. It relates to the second covariant derivative
  of $T$ (\cref{def:pet-ch2-second-covariant-derivative}) by
  $\nabla^2_{X,Y}T=(\nabla_X\nabla)_YT+\nabla_Y(\nabla_XT)$.
\end{definition}

\section{Natural Derivations}

The pointwise action of endomorphisms packages Lie and covariant derivatives
as derivations of the tensor algebra.

\subsection{Endomorphisms as Derivations}

\begin{definition}[derivation on tensors]
  \label{def:pet-ch2-tensor-derivation}
  \lean{PetersenLib.IsTensorDerivation}
  \source{petersen:page-0080}
  \leanok
  A \emph{derivation} on tensors is a type-preserving linear map
  $T\mapsto DT$ that commutes with contractions and satisfies the product rule
  $D(T_1\otimes T_2)=(DT_1)\otimes T_2+T_1\otimes(DT_2)$.
\end{definition}

\begin{definition}[the natural $\mathrm{GL}(V)$-action on $T(V)$]
  \label{def:pet-ch2-gl-action-tensor-algebra}
  \lean{PetersenLib.glActionOnTensorAlgebra}
  \source{petersen:page-0080,petersen:page-0081}
  \leanok
  For a finite-dimensional vector space $V$, let $T(V)$ be its tensor algebra,
  graded by $(s,t)$-tensors spanned by
  $v_1\otimes\cdots\otimes v_s\otimes\phi_1\otimes\cdots\otimes\phi_t$
  ($v_i\in V$, $\phi_j\in V^*$). The natural homomorphism
  $\mathrm{GL}(V)\to\mathrm{GL}(T(V))$ acts by $g\cdot\alpha=\alpha$ for
  $\alpha\in\mathbb{R}$ (a $(0,0)$-tensor); $g\cdot v=g(v)$ for $v\in V$;
  $g\cdot\phi=\phi\circ g^{-1}$ for $\phi\in V^*$; and, on general tensors,
  \[
    g\cdot(v_1\otimes\cdots\otimes v_s\otimes\phi_1\otimes\cdots\otimes\phi_t)
    = g(v_1)\otimes\cdots\otimes g(v_s)\otimes(\phi_1\circ g^{-1})\otimes\cdots\otimes(\phi_t\circ g^{-1}).
  \]
\end{definition}

\begin{definition}[induced derivation of $L\in\mathrm{End}(V)$]
  \label{def:pet-ch2-endomorphism-derivation}
  \lean{PetersenLib.endomorphismDerivation}
  \source{petersen:page-0081}
  \uses{def:pet-ch2-gl-action-tensor-algebra, def:pet-ch2-tensor-derivation}
  \leanok
  Differentiating the $\mathrm{GL}(V)$-action of
  \cref{def:pet-ch2-gl-action-tensor-algebra} gives a linear map
  $\mathrm{End}(V)\to\mathrm{End}(T(V))$, $L\mapsto LT$: on vectors $Lv=L(v)$;
  on $1$-forms $L\phi=-\phi\circ L$; and on general tensors,
  \[
    L(v_1\otimes\cdots\otimes v_s\otimes\phi_1\otimes\cdots\otimes\phi_t)
    = \sum_{i=1}^sv_1\otimes\cdots\otimes L(v_i)\otimes\cdots\otimes v_s\otimes\phi_1\otimes\cdots\otimes\phi_t
  \]
  \[
    -\sum_{j=1}^tv_1\otimes\cdots\otimes v_s\otimes\phi_1\otimes\cdots\otimes(\phi_j\circ L)\otimes\cdots\otimes\phi_t.
  \]
\end{definition}

\begin{proposition}[Lie algebra homomorphism]
  \label{prop:pet-ch2-endomorphism-lie-algebra-hom}
  \lean{PetersenLib.endomorphismDerivation_lieAlgebraHom}
  \source{petersen:page-0081}
  \uses{def:pet-ch2-endomorphism-derivation}
  \leanok
  \dcref{ch2:2.3.1}
  The linear map $\mathrm{End}(V)\to\mathrm{End}(T(V))$, $L\mapsto LT$, is a
  Lie algebra homomorphism that preserves symmetries of tensors.
\end{proposition}
\begin{proof}
  \uses{def:pet-ch2-endomorphism-derivation}
  \leanok
  \sketch
  In
  $L_1(L_2T)-L_2(L_1T)$, terms in which the maps act on different tensor
  slots cancel in pairs. Terms acting twice on one slot combine to
  $(L_1L_2-L_2L_1)T$. Thus the induced map preserves brackets. It also
  commutes with permutations of tensor slots and therefore preserves tensor
  symmetries.
\end{proof}

\begin{proposition}[$(1,1)$-tensors act as derivations]
  \label{prop:pet-ch2-endomorphism-is-derivation}
  \lean{PetersenLib.endomorphismAction_isDerivation}
  \source{petersen:page-0081,petersen:page-0082}
  \uses{def:pet-ch2-endomorphism-derivation, def:pet-ch2-tensor-derivation}
  \leanok
  \dcref{ch2:2.3.2}
  Any $(1,1)$-tensor $L$ defines a derivation on tensors
  (\cref{def:pet-ch2-tensor-derivation}).
\end{proposition}
\begin{proof}
  \uses{prop:pet-ch2-endomorphism-is-derivation}
  \leanok
  \sketch
  Linearity and the product
  rule follow from the slot formula. For contraction, in a frame write
  $T=T^i_jX_i\otimes\sigma^j$ and $L(X_i)=L^k_iX_k$. Then
  \[
    L(T)
      =T^i_jL^k_iX_k\otimes\sigma^j
       -T^i_jL^j_kX_i\otimes\sigma^k,
    \qquad
    \operatorname{tr}(L(T))=T^i_kL^k_i-T^i_kL^k_i=0.
  \]
  This equals $L(\operatorname{tr}T)$ because the induced action annihilates
  scalars. The same cancellation applies to any chosen upper-lower
  contraction.
\end{proof}

\begin{definition}[inner product on the tensor algebra]
  \label{def:pet-ch2-tensor-algebra-inner-product}
  \lean{PetersenLib.tensorAlgebraInnerProduct}
  \source{petersen:page-0082}
  \leanok
  If $V$ has an inner product with orthonormal basis $e_1,\dots,e_n$ and dual
  basis $e^1,\dots,e^n$, the inner product on $T(V)$ declares
  $e_{i_1}\otimes\cdots\otimes e_{i_p}\otimes e^{j_1}\otimes\cdots\otimes e^{j_q}$
  an orthonormal basis of the $(p,q)$-tensors.
\end{definition}

\begin{proposition}[Adjoints and skew-adjoint endomorphisms]
  \label{prop:pet-ch2-endomorphism-adjoint-type-change}
  \lean{PetersenLib.endomorphismDerivation_adjoint_skew}
  \source{petersen:page-0082}
  \uses{def:pet-ch2-endomorphism-derivation, def:pet-ch2-tensor-algebra-inner-product}
  \leanok
  \dcref{ch2:2.3.3}
  Assume $V$ has an inner product. (1) The adjoint of $L:V\to V$ extends to
  the adjoint of $L:T(V)\to T(V)$. (2) If $L\in\mathfrak{so}(V)$ (skew-adjoint),
  then $L$ commutes with type change of tensors.
\end{proposition}
\begin{proof}
  \uses{def:pet-ch2-endomorphism-derivation,def:pet-ch2-tensor-algebra-inner-product}
  \leanok
  \sketch
  In an orthonormal tensor
  basis the induced action is the sum of the matrix of $L$ over all slots;
  transposing each summand gives the induced action of $L^\dagger$. Hence
  $\langle LT_1,T_2\rangle=\langle T_1,L^\dagger T_2\rangle$. If
  $L^\dagger=-L$, the sign from moving $L$ through the metric cancels the
  sign attached to a covariant slot, so raising or lowering an index commutes
  with the induced action.
\end{proof}

\subsection{Derivatives}

\begin{proposition}[$L_U$ decomposes through $\nabla$]
  \label{prop:pet-ch2-lie-derivative-as-endomorphism-derivation}
  \lean{PetersenLib.lieDerivative_eq_covariant_minus_endomorphism}
  \source{petersen:page-0082,petersen:page-0083}
  \uses{def:pet-ch2-lie-derivative-vector-field, def:pet-ch2-covariant-derivative-tensor, def:pet-ch2-endomorphism-derivation}
  \leanok
  \dcref{ch2:2.3.4}
  Viewing $\nabla U$ as the $(1,1)$-tensor $X\mapsto\nabla_XU$, the Lie
  derivative in the direction $U$ decomposes as $L_U=\nabla_U-(\nabla U)$,
  where $(\nabla U)$ acts by the endomorphism derivation of
  \cref{def:pet-ch2-endomorphism-derivation}.
\end{proposition}
\begin{proof}
  \uses{prop:pet-ch2-lie-derivative-as-endomorphism-derivation}
  \leanok
  \sketch
  Both sides are
  tensor derivations, so it suffices to compare them on functions and vector
  fields. On functions, $(\nabla U)$ acts trivially and
  $L_Uf=\nabla_Uf$. On a vector field $Y$, torsion-freeness gives
  \[
    L_UY=[U,Y]=\nabla_UY-\nabla_YU
      =\big(\nabla_U-(\nabla U)\big)Y.
  \]
\end{proof}

\section{The Connection in Tensor Notation}

Fix local coordinates $x^1,\dots,x^n$ and use the Einstein convention.

\begin{definition}[coordinate formulas for the metric]
  \label{def:pet-ch2-metric-coordinate-formulas}
  \lean{PetersenLib.metricCoordinateFormulas}
  \source{petersen:page-0083}
  \leanok
  In local coordinates,
  \[
    g=g_{ij}dx^i dx^j,\qquad
    g(X,Y)=g_{ij}X^iY^j,\qquad
    \theta_X=g_{ij}X^i dx^j,
  \]
  where $\theta_X$ is the dual form of \cref{def:pet-ch2-dual-one-form}.
  The inverse coefficients satisfy $g^{ik}g_{kj}=\delta^i_j$.
\end{definition}

\begin{definition}[vector field dual to a $1$-form]
  \label{def:pet-ch2-vector-dual-to-form}
  \lean{PetersenLib.vectorDualToForm}
  \source{petersen:page-0083}
  \uses{def:pet-ch2-metric-coordinate-formulas}
  \leanok
  The vector field dual to
  $\omega=\omega_jdx^j$ is
  \[
    \omega^\sharp=g^{ij}\omega_j\partial_i,
  \]
  since $g_{ij}(\omega^\sharp)^i=\omega_j$.
\end{definition}

\begin{remark}[gradient in coordinates]
  \label{rem:pet-ch2-gradient-coordinate-formula}
  \lean{PetersenLib.gradient_coordinate_formula}
  \source{petersen:page-0084}
  \uses{def:pet-ch2-directional-derivative, def:pet-ch2-vector-dual-to-form}
  \leanok
  Applying \cref{def:pet-ch2-vector-dual-to-form} to $\omega=df=\partial_jf\,dx^j$
  gives $\nabla f=g^{ij}\partial_jf\,\partial_i$, matching
  \cref{def:pet-ch2-directional-derivative}.
\end{remark}

\begin{definition}[Christoffel symbols of the second kind]
  \label{def:pet-ch2-christoffel-symbols-second-kind}
  \lean{PetersenLib.christoffelSymbolsSecondKind}
  \source{petersen:page-0084}
  \uses{def:pet-ch2-covariant-derivative-tensor, def:pet-ch2-metric-coordinate-formulas}
  \leanok
  In local coordinates, expand
  $\nabla_YX=Y^i\big(\partial_i(X^j)\big)\partial_j+Y^iX^j\Gamma^k_{ij}\partial_k$, where
  $\Gamma^k_{ij}$, the \emph{Christoffel symbols of the second kind}, are
  defined by $\nabla_{\partial_i}\partial_j=\Gamma^k_{ij}\partial_k$.
\end{definition}

\begin{proposition}[$\Gamma^k_{ij}$ in terms of the metric]
  \label{prop:pet-ch2-christoffel-formula}
  \lean{PetersenLib.christoffelSymbols_metric_formula}
  \source{petersen:page-0084,petersen:page-0085}
  \uses{def:pet-ch2-christoffel-symbols-second-kind, rem:pet-ch2-koszul-formula}
  \leanok
  $\displaystyle\Gamma^k_{ij}=\tfrac12g^{k\ell}\big(\partial_ig_{\ell j}+\partial_jg_{\ell i}-\partial_\ell g_{ij}\big)$.
\end{proposition}
\begin{proof}
  \uses{prop:pet-ch2-christoffel-formula}
  \leanok
  \sketch
  Apply Koszul's formula to the
  commuting coordinate fields $\partial_i,\partial_j,\partial_\ell$:
  \[
    2\Gamma^k_{ij}g_{k\ell}
      =\partial_i g_{j\ell}+\partial_j g_{i\ell}
       -\partial_\ell g_{ij}.
  \]
  Contracting with $g^{\ell m}$ gives the asserted formula.
\end{proof}

\begin{definition}[Christoffel symbols of the first kind]
  \label{def:pet-ch2-christoffel-symbols-first-kind}
  \lean{PetersenLib.christoffelSymbolsFirstKind}
  \source{petersen:page-0085}
  \uses{def:pet-ch2-christoffel-symbols-second-kind}
  \leanok
  $\Gamma_{ij,k}:=\tfrac12(\partial_jg_{ik}+\partial_ig_{jk}-\partial_kg_{ij})=g(\nabla_{\partial_i}\partial_j,\partial_k)$;
  classically also written $[ij,k]$, with $\Gamma^k_{ij}=\begin{Bmatrix}k\\i,j\end{Bmatrix}$.
\end{definition}

\begin{remark}[Christoffel symbols are not tensorial]
  \label{rem:pet-ch2-christoffel-non-tensorial}
  \lean{PetersenLib.christoffelSymbols_not_tensorial}
  \source{petersen:page-0085}
  \uses{def:pet-ch2-christoffel-symbols-first-kind, rem:pet-ch2-ex-22}
  \leanok
  Christoffel symbols are
  coordinate-dependent. They vanish in Cartesian coordinates on
  $\mathbb{R}^2$, whereas for the polar parametrization
  $u(r,\theta)=(r\cos\theta,r\sin\theta)$,
  \[
    g=dr^2+r^2d\theta^2,\qquad
    \Gamma_{\theta\theta,r}=-\tfrac12\partial_r(r^2)=-r.
  \]
  \sketch
  The formalized comparison uses the induced metric
  $g_x(v,w)=\langle Du_xv,Du_xw\rangle$ from
  \cref{rem:pet-ch2-ex-22}. The Cartesian chart has zero first-kind symbols,
  while the displayed polar symbol is nonzero away from the origin. Since
  coordinate components of a tensor cannot vanish in only one chart, the
  Christoffel symbols are not tensor components.
\end{remark}

\begin{remark}[normal coordinates and the Christoffel symbols]
  \label{rem:pet-ch2-christoffel-normal-coordinates}
  \lean{PetersenLib.christoffelSymbols_vanish_normalCoordinates}
  \source{petersen:page-0085}
  \uses{def:pet-ch2-normal-coordinates-frame, def:pet-ch2-christoffel-symbols-second-kind}
  \leanok
  In coordinates normal at $p$ (\cref{def:pet-ch2-normal-coordinates-frame}),
  $g_{ij}|_p=\delta_{ij}$ and $\Gamma^k_{ij}|_p=0$; consequently, at $p$, the
  covariant derivative is computed exactly as in Euclidean space:
  $(\nabla_YX)|_p=\big(Y^i\partial_i(X^j)\big)|_p\,\partial_j|_p$.
\end{remark}

\begin{proposition}[torsion-free and metric properties in coordinates]
  \label{prop:pet-ch2-christoffel-symmetry-metric-property}
  \lean{PetersenLib.christoffel_symmetric_metric_property}
  \source{petersen:page-0086}
  \uses{def:pet-ch2-christoffel-symbols-first-kind, def:pet-ch2-christoffel-symbols-second-kind}
  \leanok
  Torsion-freeness (\cref{def:pet-ch2-riemannian-connection}(3)) is equivalent
  to $\Gamma^k_{ij}=\Gamma^k_{ji}$, and the metric property
  (\cref{def:pet-ch2-riemannian-connection}(4)) becomes
  $\partial_kg_{ij}=\Gamma_{ki,j}+\Gamma_{kj,i}$, with the first- and
  second-kind symbols as in
  \cref{def:pet-ch2-christoffel-symbols-first-kind} and
  \cref{def:pet-ch2-christoffel-symbols-second-kind}.
\end{proposition}
\begin{proof}
  \uses{prop:pet-ch2-christoffel-symmetry-metric-property}
  \leanok
  \sketch
  Since
  $[\partial_i,\partial_j]=0$, torsion-freeness is equivalent to
  $\nabla_{\partial_i}\partial_j=\nabla_{\partial_j}\partial_i$, hence to
  $\Gamma^k_{ij}=\Gamma^k_{ji}$. Metric compatibility applied to
  $g(\partial_i,\partial_j)$ gives
  \[
    \partial_k g_{ij}
      =g(\nabla_{\partial_k}\partial_i,\partial_j)
       +g(\partial_i,\nabla_{\partial_k}\partial_j)
      =\Gamma_{ki,j}+\Gamma_{kj,i}.
  \]
\end{proof}

\begin{proposition}[Hessian in Christoffel-symbol coordinates]
  \label{prop:pet-ch2-hessian-coordinate-formula}
  \lean{PetersenLib.hessian_coordinate_formula}
  \source{petersen:page-0086,petersen:page-0087}
  \uses{def:pet-ch2-hessian-lie-derivative, def:pet-ch2-christoffel-symbols-second-kind, prop:pet-ch2-hessian-covariant-derivative}
  \leanok
  $\operatorname{Hess}f(\partial_i,\partial_j)=\partial_i\partial_jf-\Gamma^k_{ij}\partial_kf$,
  the $\partial_i$ being the coordinate frame.
\end{proposition}
\begin{proof}
  \leanok
  \uses{prop:pet-ch2-hessian-coordinate-formula}
  \sketch
  By
  \cref{prop:pet-ch2-hessian-covariant-derivative},
  \[
    \operatorname{Hess}f(X,Y)
      =D_XD_Yf-D_{\nabla_XY}f.
  \]
  Substituting $X=\partial_i$, $Y=\partial_j$, and the second-kind definition
  \cref{def:pet-ch2-christoffel-symbols-second-kind},
  $\nabla_{\partial_i}\partial_j=\Gamma^k_{ij}\partial_k$ gives the claimed
  coordinate expression. Substitution from
  \cref{prop:pet-ch2-christoffel-formula} rewrites the correction in terms of
  the first derivatives of the metric.
\end{proof}

\begin{definition}[index notation for the covariant derivative]
  \label{def:pet-ch2-covariant-derivative-index-notation}
  \lean{PetersenLib.covariantDerivativeIndexNotation}
  \source{petersen:page-0087,petersen:page-0088}
  \uses{def:pet-ch2-covariant-derivative-tensor}
  \leanok
  In a frame $E_i$
  with coframe $\sigma^j$, write
  \[
    S=S^i_{j_1\cdots j_k}
      E_i\otimes\sigma^{j_1}\otimes\cdots\otimes\sigma^{j_k}.
  \]
  The coefficients of $\nabla S$ are obtained by applying the product rule to
  the scalar coefficient and every frame or coframe factor. We use
  \[
    \nabla_jS:=\nabla_{E_j}S,\qquad
    \nabla_jS^i_{j_1\cdots j_k}
      :=(\nabla S)^i_{j_1\cdots j_kj},
  \]
  so the differentiation index is written first in operator notation.
\end{definition}

\section{Exercises}

\begin{remark}[Uniqueness of the Euclidean connection]
  \label{rem:pet-ch2-ex-1}
  \lean{PetersenLib.exercise2_5_1}
  \source{petersen:page-0088}
  \uses{def:pet-ch2-covariant-derivative-euclidean, def:pet-ch2-affine-connection}
  \leanok
  \dcref{ch2:2.5.1}
  The Cartesian connection is the unique affine
  connection on Euclidean space for which every constant vector field is
  parallel.
  \sketch
  For $X=X^i\partial_i$, the Leibniz rule and
  $\nabla\partial_i=0$ force
  $\nabla_YX=dX^i(Y)\partial_i$, which is
  \cref{def:pet-ch2-covariant-derivative-euclidean}.
\end{remark}

\begin{remark}[Regularity needed for skew-symmetry and Jacobi]
  \label{rem:pet-ch2-ex-2}
  \lean{PetersenLib.exercise2_5_2}
  \source{petersen:page-0088}
  \uses{prop:pet-ch2-lie-derivative-bracket}
  \leanok
  \dcref{ch2:2.5.2}
  Under the smooth identification in
  \cref{prop:pet-ch2-lie-derivative-bracket}, Petersen's flow-defined $L_XY$ need not be
  skew-symmetric for merely $C^1$ fields; its identification with the coordinate
  commutator requires the relevant second derivatives. For $C^2$ fields,
  \[
    [X,[Y,W]]=[[X,Y],W]+[Y,[X,W]],
  \]
  equivalently the cyclic Jacobi sum vanishes.
  The formalization uses Mathlib's coordinate bracket
  $\texttt{VectorField.mlieBracket}$, whose skew-symmetry is unconditional;
  hence the requested $C^1$ counterexample is outside its scope. The $C^2$
  Jacobi identity is formalized.
\end{remark}

\begin{remark}[The torsion tensor]
  \label{rem:pet-ch2-ex-3}
  \lean{PetersenLib.exercise2_5_3}
  \source{petersen:page-0088}
  \uses{def:pet-ch2-affine-connection}
  \leanok
  \dcref{ch2:2.5.3}
  For any affine connection,
  \[
    T(X,Y)=\nabla_XY-\nabla_YX-[X,Y]
  \]
  is a $(2,1)$-tensor.
  \sketch
  Additivity is immediate. In each argument, the derivative terms generated
  by multiplying a field by a function cancel the corresponding Lie-bracket
  Leibniz term, leaving $C^\infty(M)$-linearity.
\end{remark}

\begin{remark}[Extending a tangent to a field]
  \label{rem:pet-ch2-ex-4}
  \lean{PetersenLib.exercise2_5_4}
  \source{petersen:page-0088}
  \leanok
  \dcref{ch2:2.5.4}
  If $c:I\to M$ satisfies
  $\dot c(t_0)\ne0$, then some smooth vector field $X$ obeys
  $X_{c(t)}=\dot c(t)$ for all $t$ near $t_0$.
  \sketch
  In a chart, choose a linear functional nonzero on $\dot c(t_0)$ and use the
  inverse function theorem to rectify the curve locally. Its coordinate
  velocity then extends to a smooth field on a neighborhood; pull it back
  through the chart, multiply by a cutoff, and extend by zero. This complete
  local construction is formalized.
\end{remark}

\begin{remark}[Product rules for divergence, Laplacian, Hessian]
  \label{rem:pet-ch2-ex-5}
  \lean{PetersenLib.exercise2_5_5}
  \source{petersen:page-0088}
  \uses{def:pet-ch2-divergence-lie-derivative, def:pet-ch2-laplacian, def:pet-ch2-hessian-lie-derivative}
  \leanok
  \dcref{ch2:2.5.5}
  For functions $f,h$ and a vector field $X$,
  \[
  \begin{aligned}
    \operatorname{div}(fX)
      &=D_Xf+f\operatorname{div}X,\\
    \Delta(fh)
      &=h\Delta f+f\Delta h+2g(\nabla f,\nabla h),\\
    \operatorname{Hess}(fh)
      &=h\operatorname{Hess}f+f\operatorname{Hess}h
        +df\otimes dh+dh\otimes df.
  \end{aligned}
  \]
  \sketch
  These follow from the Leibniz rules for $L_X$, $\nabla$, and the gradient.
\end{remark}

\begin{remark}[Chain rule for Laplacian and Hessian]
  \label{rem:pet-ch2-ex-6}
  \lean{PetersenLib.exercise2_5_6}
  \source{petersen:page-0089}
  \uses{def:pet-ch2-laplacian, def:pet-ch2-hessian-lie-derivative}
  \leanok
  \dcref{ch2:2.5.6}
  For $f:M\to\mathbb{R}$ and
  $\phi:\mathbb{R}\to\mathbb{R}$,
  \[
  \begin{aligned}
    \Delta(\phi\circ f)
      &=\dot\phi(f)\Delta f+\ddot\phi(f)|df|^2,\\
    \operatorname{Hess}(\phi\circ f)
      &=\dot\phi(f)\operatorname{Hess}f+\ddot\phi(f)\,df^2.
  \end{aligned}
  \]
\end{remark}

\begin{remark}[$d\theta_X$ via the connection]
  \label{rem:pet-ch2-ex-7}
  \lean{PetersenLib.exercise2_5_7}
  \source{petersen:page-0089}
  \uses{def:pet-ch2-dual-one-form, def:pet-ch2-covariant-derivative-tensor}
  \leanok
  \dcref{ch2:2.5.7}
  If $\theta_X=g(X,\cdot)$, then
  \[
    d\theta_X(Y,Z)
      =g(\nabla_YX,Z)-g(Y,\nabla_ZX).
  \]
  \sketch
  This follows by expanding $d\theta_X$, applying metric compatibility, and
  replacing $[Y,Z]$ by $\nabla_YZ-\nabla_ZY$.
\end{remark}

\begin{remark}[Derivatives of the inverse metric coefficients]
  \label{rem:pet-ch2-ex-8}
  \lean{PetersenLib.exercise2_5_8}
  \source{petersen:page-0089}
  \uses{def:pet-ch2-metric-coordinate-formulas, def:pet-ch2-christoffel-symbols-second-kind}
  \leanok
  \dcref{ch2:2.5.8}
  The inverse metric satisfies
  \[
    \partial_kg^{ij}
      =-g^{i\ell}\Gamma^j_{k\ell}
       -g^{j\ell}\Gamma^i_{k\ell}.
  \]
  \sketch
  Differentiate $g^{i\ell}g_{\ell j}=\delta^i_j$ and use metric
  compatibility in coordinates.
\end{remark}

\begin{remark}[Covariant differentiation commutes with contraction and type change]
  \label{rem:pet-ch2-ex-9}
  \lean{PetersenLib.exercise2_5_9}
  \source{petersen:page-0089}
  \uses{def:pet-ch2-covariant-derivative-tensor}
  \leanok
  \dcref{ch2:2.5.9}
  For a vector field $X$ and a $(1,1)$-tensor
  $S$:
  \begin{enumerate}
    \item $\operatorname{tr}(\nabla_XS)=\nabla_X(\operatorname{tr}S)$;
    \item if $T(Y,Z)=g(S(Y),Z)$, then
      $(\nabla_XT)(Y,Z)=g((\nabla_XS)(Y),Z)$;
    \item contraction commutes with covariant differentiation;
    \item type change commutes with covariant differentiation.
  \end{enumerate}
  Parts (1)--(2) are formalized; the general tensor operations required for
  (3)--(4) are not.
  \sketch
  Metric compatibility expands the derivative of $g(S(Y),Z)$, and the
  corrections in the $Y$- and $Z$-slots leave
  $g(\nabla_X(S(Y))-S(\nabla_XY),Z)$. For the trace, differentiate
  $\sum_i g(E_i,S(E_i))$ in an orthonormal frame. The remaining frame terms
  cancel using
  $g(\nabla_XE_i,E_j)+g(E_i,\nabla_XE_j)=0$.
\end{remark}

\begin{remark}[Lie differentiation commutes with contraction]
  \label{rem:pet-ch2-ex-10}
  \lean{PetersenLib.exercise2_5_10}
  \source{petersen:page-0089}
  \uses{def:pet-ch2-lie-derivative-tensor}
  \leanok
  \dcref{ch2:2.5.10}
  For a vector field $X$ and a $(1,1)$-tensor
  $S$:
  \begin{enumerate}
    \item $\operatorname{tr}(L_XS)=L_X(\operatorname{tr}S)$;
    \item for $T(Y,Z)=g(S(Y),Z)$,
      \[
        (L_XT)(Y,Z)=(L_Xg)(S(Y),Z)+g((L_XS)(Y),Z);
      \]
    \item contraction commutes with Lie differentiation.
  \end{enumerate}
  Parts (1)--(2) are formalized; (3) awaits the general mixed-tensor
  contraction layer.
  \sketch
  Expanding the three Lie derivatives proves (2) by cancellation. For (1),
  torsion-freeness gives
  \[
    L_XS=\nabla_XS-(\nabla X)\circ S+S\circ(\nabla X).
  \]
  Taking traces, cyclicity cancels the last two terms, and
  \cref{rem:pet-ch2-ex-9} gives the result.
\end{remark}

\begin{remark}[Local gradient fields via self-adjointness]
  \label{rem:pet-ch2-ex-11}
  \lean{PetersenLib.exercise2_5_11}
  \source{petersen:page-0089}
  \uses{def:pet-ch2-covariant-derivative-tensor}
  \leanok
  \dcref{ch2:2.5.11}
  A vector field $X$ is locally a gradient
  exactly when $Z\mapsto\nabla_ZX$ is self-adjoint.
  The forward implication is formalized; the converse is the local
  Poincar\'e lemma and is deferred.
  \sketch
  If $X=\nabla f$, then
  \[
    g(\nabla_U\nabla f,V)=g(\nabla_V\nabla f,U)
  \]
  by symmetry of $\nabla df$ from
  \cref{prop:pet-ch2-hessian-covariant-derivative}.
\end{remark}

\begin{remark}[Isometries and pushforward of the connection]
  \label{rem:pet-ch2-ex-12}
  \lean{PetersenLib.exercise2_5_12}
  \source{petersen:page-0089,petersen:page-0090}
  \uses{def:pet-ch2-covariant-derivative-tensor, thm:pet-ch2-fundamental-theorem}
  \leanok
  \dcref{ch2:2.5.12}
  For a diffeomorphism $F$, define
  $(F_*X)_p=DF_{F^{-1}(p)}X_{F^{-1}(p)}$. If $F:(M,g)\to(M',g')$ is a
  Riemannian isometry, then
  \[
    F_*(\nabla^g_YX)=\nabla^{g'}_{F_*Y}(F_*X).
  \]
  Consequently every Euclidean isometry is $F(x)=Ox+b$ with
  $O\in O(n)$ and $b\in\mathbb{R}^n$. Both assertions are formalized.
  \sketch
  The metric pairing, directional derivatives, and Lie brackets in Koszul's
  formula are natural under $F$. Uniqueness in
  \cref{thm:pet-ch2-fundamental-theorem} therefore proves the connection
  identity. The Euclidean classification follows by sending constant fields
  to constant fields, equivalently from the Chapter~1 Euclidean isometry
  theorem.
\end{remark}

\begin{remark}[Affine vector fields]
  \label{rem:pet-ch2-ex-13}
  \lean{PetersenLib.exercise2_5_13}
  \source{petersen:page-0090}
  \uses{def:pet-ch2-lie-derivative-connection, def:pet-ch2-killing-field, def:pet-ch3-curvature-tensor, rem:pet-ch3-ricci-identity-derivation, prop:pet-ch3-curvature-symmetries, rem:pet-ch3-ex-19, def:pet-ch2-covariant-derivative-euclidean, rem:pet-ch2-killing-field-euclidean, ex:pet-ch3-euclidean-flat}
  \leanok
  \dcref{ch2:2.5.13}
  A vector field is \emph{affine} when
  $L_X\nabla=0$. Every Killing field is affine. The converse fails on
  $\mathbb{R}^n$: the radial field $X(x)=x$ is affine but not Killing. Both
  statements are formalized.
  \sketch
  For Killing $X$, $\nabla X$ is skew-adjoint. Differentiation, the Ricci
  identity \cref{rem:pet-ch3-ricci-identity-derivation}, and the curvature
  symmetries of \cref{prop:pet-ch3-curvature-symmetries} give Kostant's formula
  $\nabla^2_{U,V}X=-R(X,U)V$. Together with
  \[
    (L_X\nabla)(U,V)=R(X,U)V+\nabla^2_{U,V}X
  \]
  from \cref{rem:pet-ch3-ex-19}, this gives $L_X\nabla=0$. For $X(x)=x$,
  $\nabla_VX=V$ and $\nabla^2X=0$; Euclidean flatness makes it affine, while
  $(L_Xg)(v,v)=2|v|^2$.
\end{remark}

\begin{remark}[The left-invariant connection on a Lie group]
  \label{rem:pet-ch2-ex-14}
  \lean{PetersenLib.exercise2_5_14}
  \source{petersen:page-0090}
  \uses{def:pet-ch2-affine-connection}
  \leanok
  \dcref{ch2:2.5.14}
  Every Lie group $G$ has a unique affine
  connection for which all left-invariant vector fields are parallel. This
  connection is torsion free exactly when the Lie algebra is Abelian.
  \sketch
  For a basis $e_i$ of $\mathfrak g=T_eG$, let $E_i$ be the corresponding
  left-invariant frame. Writing $X=X^iE_i$, define
  \[
    \nabla_vX=\sum_i dX^i(v)E_i.
  \]
  This is an affine connection and kills every left-invariant field.
  The Leibniz rule in the same frame proves uniqueness. For left-invariant
  $X,Y$, its torsion is $T(X,Y)=-[X,Y]$, proving the final equivalence. The
  entire construction is formalized.
\end{remark}

\begin{remark}[Chain rule for the Hessian]
  \label{rem:pet-ch2-ex-15}
  \lean{PetersenLib.exercise2_5_15}
  \source{petersen:page-0090}
  \uses{def:pet-ch2-hessian-lie-derivative}
  \leanok
  \dcref{ch2:2.5.15}
  For $\phi:\mathbb{R}\to\mathbb{R}$,
  \[
    \operatorname{Hess}(\phi\circ f)
      =\phi''(f)\,df^2+\phi'(f)\operatorname{Hess}f.
  \]
\end{remark}

\begin{remark}[Derivations are $L_X+L$]
  \label{rem:pet-ch2-ex-16}
  \lean{PetersenLib.exercise2_5_16}
  \source{petersen:page-0090}
  \uses{def:pet-ch2-lie-derivative-tensor, prop:pet-ch2-lie-derivative-scaling}
  \leanok
  \dcref{ch2:2.5.16}
  Let $X$ be a vector field and $L$ a
  $(1,1)$-tensor:
  \begin{enumerate}
    \item $L_X+L$ is a tensor derivation in the sense of
      \cref{def:pet-ch2-tensor-derivation};
    \item every derivation has this form with $X$ unique;
    \item a derivation is determined on functions and vector fields;
    \item
      \[
        L_{fX}=fL_X-X\otimes df,\qquad
        (X\otimes df)(Y)=df(Y)X.
      \]
  \end{enumerate}
  Part (4) is formalized for $(0,k)$-tensors: the rank-one derivation acts by
  inserting $df(Y_i)X$ in each covariant slot, and the formula is
  \cref{prop:pet-ch2-lie-derivative-scaling}. Parts (1)--(3) require a
  manifold-level derivation across all tensor types and are deferred.
\end{remark}

\begin{remark}[Constant-length fields are orthogonal to their covariant derivative]
  \label{rem:pet-ch2-ex-17}
  \lean{PetersenLib.exercise2_5_17}
  \source{petersen:page-0090}
  \uses{def:pet-ch2-covariant-derivative-tensor}
  \leanok
  \dcref{ch2:2.5.17}
  If $X$ has constant length, then
  $g(\nabla_ZX,X)=0$ for every $Z$.
  \sketch
  Metric compatibility gives
  $D_Zg(X,X)=2g(\nabla_ZX,X)$, and the left side vanishes.
\end{remark}

\begin{remark}[Lie and covariant derivatives agree where $T$ vanishes]
  \label{rem:pet-ch2-ex-18}
  \lean{PetersenLib.exercise2_5_18}
  \source{petersen:page-0090}
  \uses{def:pet-ch2-lie-derivative-tensor, def:pet-ch2-covariant-derivative-tensor, def:pet-ch2-hessian-lie-derivative}
  \leanok
  \dcref{ch2:2.5.18}
  If a tensor field $T$ vanishes at $p$, then
  \[
    (L_XT)_p=(\nabla_XT)_p.
  \]
  \sketch
  Indeed, \cref{prop:pet-ch2-lie-derivative-as-endomorphism-derivation} gives
  $L_XT=\nabla_XT-(\nabla X)T$, and the algebraic final term vanishes with
  $T_p$. Applied at a critical point to the $(1,1)$ Hessian, this identifies
  its metric-independent first variation.
\end{remark}

\begin{remark}[Existence of a normal frame]
  \label{rem:pet-ch2-ex-19}
  \lean{PetersenLib.exercise2_5_19}
  \source{petersen:page-0090}
  \uses{def:pet-ch2-normal-coordinates-frame}
  \leanok
  \dcref{ch2:2.5.19}
  Given $p\in M$ and an orthonormal basis
  $e_1,\dots,e_n$ of $T_pM$, there is a frame $E_i$ with
  $E_i(p)=e_i$ and $(\nabla E_i)_p=0$, hence is normal at $p$ in the sense of
  \cref{def:pet-ch2-normal-coordinates-frame}.
  \sketch
  Extend $e_i$ to fields $F_i$ and set $E_i=c_{im}F_m$, where
  \[
    c_{im}(p)=\delta_{im},\qquad
    dc_{im}|_p(v)=-g(\nabla_vF_i,e_m).
  \]
  Functions with these prescribed $1$-jets exist by a chart and cutoff. The
  Leibniz rule and the orthonormal expansion at $p$ give
  $\nabla_vE_i|_p=0$. This normal-at-$p$ frame is formalized; orthonormality
  throughout a neighborhood would require an additional orthogonal gauge and
  is not separately formalized.
\end{remark}

\begin{remark}[Existence of normal coordinates]
  \label{rem:pet-ch2-ex-20}
  \lean{PetersenLib.exercise2_5_20}
  \source{petersen:page-0090,petersen:page-0091}
  \uses{def:pet-ch2-normal-coordinates-frame, rem:pet-ch2-ex-21, def:pet-ch2-christoffel-symbols-first-kind}
  \leanok
  \dcref{ch2:2.5.20}
  At every $p\in(M,g)$ there are coordinates
  $x^i$ with $\partial_i|_p=e_i$,
  $(\nabla\partial_i)|_p=0$, $g_{ij}(p)=\delta_{ij}$, and
  $\partial_kg_{ij}(p)=0$.
  \sketch
  After the linear orthonormalization supplied by
  \cref{rem:pet-ch2-ex-19}, write the coordinate metric as $B$ with
  $B(0)=\langle\cdot,\cdot\rangle$. Let $q$ be determined by
  $\langle q(v,w),z\rangle=\Gamma(v,w,z)$ and set
  $\varphi(x)=x-\tfrac12q(x,x)$. Then $D\varphi_0=I$ and the transformation
  law of \cref{rem:pet-ch2-ex-21} gives $\widetilde\Gamma_0=0$. Since
  \[
    \partial_z\widetilde g(v,w)
      =\widetilde\Gamma(z,v,w)+\widetilde\Gamma(z,w,v),
  \]
  all first metric derivatives vanish at $0$. The quadratic correction is
  formalized under the stated orthonormalized hypothesis.
\end{remark}

\begin{remark}[Transformation of Christoffel symbols]
  \label{rem:pet-ch2-ex-21}
  \lean{PetersenLib.exercise2_5_21}
  \source{petersen:page-0091}
  \uses{def:pet-ch2-christoffel-symbols-second-kind, def:pet-ch2-christoffel-symbols-first-kind}
  \leanok
  \dcref{ch2:2.5.21}
  If $\varphi:\tilde x\mapsto x$ is a coordinate
  change, then
  \[
    \widetilde\Gamma^k_{ij}
      =\frac{\partial^2x^s}
              {\partial\tilde x^i\partial\tilde x^j}
         \frac{\partial\tilde x^k}{\partial x^s}
       +\frac{\partial x^s}{\partial\tilde x^i}
        \frac{\partial x^t}{\partial\tilde x^j}
        \frac{\partial\tilde x^k}{\partial x^\ell}
        \Gamma^\ell_{st},
  \]
  equivalently
  \[
    \frac{\partial^2x^r}
         {\partial\tilde x^i\partial\tilde x^j}
      =\Gamma^k_{ij}\frac{\partial x^r}{\partial\tilde x^k}
       -\frac{\partial x^r}{\partial\tilde x^k}
        \frac{\partial x^s}{\partial\tilde x^\ell}
        \widetilde\Gamma^k_{ij}
  \]
  with the indicated contractions, and
  \[
    \widetilde\Gamma_{ij,k}
      =\frac{\partial^2x^s}
              {\partial\tilde x^i\partial\tilde x^j}
         \frac{\partial x^\ell}{\partial\tilde x^k}g_{s\ell}
       +\frac{\partial x^s}{\partial\tilde x^i}
        \frac{\partial x^t}{\partial\tilde x^j}
        \frac{\partial x^\ell}{\partial\tilde x^k}\Gamma_{st,\ell}.
  \]
  \sketch
  The formalization proves the first-kind identity for the pullback metric
  $B(\varphi x)(D\varphi\,\cdot,D\varphi\,\cdot)$:
  \[
    \widetilde\Gamma(v,w,z)
      =B(D^2\varphi(v,w),D\varphi z)
       +\Gamma(D\varphi v,D\varphi w,D\varphi z).
  \]
  It follows by differentiating the pullback metric and cancelling terms using
  symmetry of $B$ and $D^2\varphi$. Raising the final index yields the
  second-kind and isolated-second-derivative forms. Taking $B$ to be the
  constant Euclidean inner product recovers \cref{rem:pet-ch2-ex-22}.
\end{remark}

\begin{remark}[Christoffel symbols of an induced metric]
  \label{rem:pet-ch2-ex-22}
  \lean{PetersenLib.exercise2_5_22}
  \source{petersen:page-0091}
  \uses{def:pet-ch2-metric-coordinate-formulas, def:pet-ch2-christoffel-symbols-first-kind}
  \leanok
  \dcref{ch2:2.5.22}
  For a parametrized submanifold
  $u:\mathbb{R}^n\to\mathbb{R}^{n+m}$ with the induced metric,
  \[
    g_{ij}=\sum_su^s_i u^s_j,\qquad
    \Gamma_{ij,k}=\sum_su^s_{ij}u^s_k.
  \]
  \sketch
  In invariant coordinate notation,
  $g_x(v,w)=\langle Du_xv,Du_xw\rangle$. Differentiating the three terms in
  the first-kind Christoffel formula and using symmetry of $D^2u$ cancels four
  terms and leaves
  \[
    \Gamma_x(v,w,z)=\langle D^2u_x(v,w),Du_xz\rangle.
  \]
  This coordinate statement is fully formalized.
\end{remark}

\begin{remark}[Volume form and Laplacian in coordinates]
  \label{rem:pet-ch2-ex-23}
  \lean{PetersenLib.exercise2_5_23}
  \source{petersen:page-0091,petersen:page-0092}
  \uses{def:pet-ch1-volume-form-oriented, def:pet-ch2-metric-coordinate-formulas, def:pet-ch2-laplacian}
  \leanok
  \dcref{ch2:2.5.23}
  On an oriented Riemannian $n$-manifold:
  \begin{enumerate}
    \item for a positively oriented tuple,
      \[
        \mathrm{vol}(v_1,\dots,v_n)
          =\sqrt{\det[g(v_i,v_j)]};
      \]
    \item in positive coordinates,
      $\mathrm{vol}=\sqrt{\det(g_{ij})}\,
        dx^1\wedge\cdots\wedge dx^n$;
    \item
      \[
        \Delta u
          =\frac1{\sqrt{\det(g_{ij})}}
             \partial_k\!\left(
               \sqrt{\det(g_{ij})}\,g^{k\ell}\partial_\ell u\right),
      \]
      and in normal coordinates
      $\Delta f(p)=\sum_i\partial_i^2f(p)$.
  \end{enumerate}
  \sketch
  Part (1) is formalized: in an oriented orthonormal basis, the volume is the
  determinant of the coordinate matrix $A$, while the Gram matrix is
  $AA^{\mathsf T}$; positivity selects the positive square root. Parts
  (2)--(3) require the Riemannian volume measure and manifold integration,
  which are deferred.
\end{remark}

\begin{remark}[Covariant derivative components of $(0,2)$- and $(1,1)$-tensors]
  \label{rem:pet-ch2-ex-24}
  \lean{PetersenLib.exercise2_5_24}
  \source{petersen:page-0092}
  \uses{def:pet-ch2-christoffel-symbols-second-kind, def:pet-ch2-covariant-derivative-index-notation}
  \leanok
  \dcref{ch2:2.5.24}
  If
  $T_{k\ell}=T(\partial_k,\partial_\ell)$, then
  \[
    (\nabla T)_{jk\ell}
      =\partial_jT_{k\ell}
       -\Gamma^m_{jk}T_{m\ell}
       -\Gamma^m_{j\ell}T_{km}.
  \]
  For a $(1,1)$-tensor,
  \[
    (\nabla T)^k_{ji}
      =\partial_jT^k_i-\Gamma^\ell_{ji}T^k_\ell
       +\Gamma^k_{j\ell}T^\ell_i.
  \]
  The $(0,2)$ formula is formalized. The mixed-tensor representation needed
  for the $(1,1)$ formula is deferred, as in \cref{rem:pet-ch2-ex-12}.
\end{remark}

\begin{remark}[Second fundamental form for isometric immersions]
  \label{rem:pet-ch2-ex-25}
  \lean{PetersenLib.exercise2_5_25}
  \source{petersen:page-0092}
  \uses{def:pet-ch2-covariant-derivative-tensor, def:pet-ch2-affine-connection, thm:pet-ch2-fundamental-theorem}
  \leanok
  \dcref{ch2:2.5.25}
  For an isometric immersion
  $F:(M,g_M)\hookrightarrow(\widetilde M,g_{\widetilde M})$, the tangential
  component of $\nabla^{\widetilde M}_XY$ is $\nabla^M_XY$. Its normal
  component
  \[
    \mathrm{I\!I}(X,Y)
      =\nabla^{\widetilde M}_XY-\nabla^M_XY
  \]
  is a symmetric $(2,1)$-tensor.
  \sketch
  The formalized algebraic part treats the difference of two torsion-free
  connections. Their Leibniz terms cancel, so the difference is tensorial;
  subtracting their torsion identities makes it symmetric, complementing the
  torsion construction in \cref{rem:pet-ch2-ex-3}. Identifying this
  difference with the normal projection for an immersion requires the
  tangent--normal splitting and is deferred.
\end{remark}

\begin{remark}[Normal connection along an immersed submanifold]
  \label{rem:pet-ch2-ex-26}
  \lean{PetersenLib.exercise2_5_26}
  \source{petersen:page-0092,petersen:page-0093}
  \uses{def:pet-ch2-covariant-derivative-tensor, rem:pet-ch2-ex-25}
  \leanok
  \dcref{ch2:2.5.26}
  Along the immersion of
  \cref{rem:pet-ch2-ex-25}, decompose, for tangent $X,Y$ and normal $V$,
  \[
    \nabla^{\widetilde M}_XV=\nabla^\perp_XV+T_XV,
    \qquad
    g_{\widetilde M}(T_XY,V)=-g_M(Y,T_XV).
  \]
  The normal connection is linear and a derivation in $V$ and tensorial in
  $X$.
  \sketch
  The displayed Weingarten identity is formalized: differentiate
  $g(Y,V)=0$ and apply metric compatibility, as in
  \cref{rem:pet-ch2-ex-17}, to obtain
  \[
    g(\nabla^{\widetilde M}_XY,V)
      =-g(Y,\nabla^{\widetilde M}_XV).
  \]
  The normal-bundle decomposition and the remaining connection properties
  await the immersion projection infrastructure.
\end{remark}

\begin{remark}[Divergence theorem and the maximum principle]
  \label{rem:pet-ch2-ex-27}
  \lean{PetersenLib.exercise2_5_27}
  \source{petersen:page-0093}
  \uses{def:pet-ch2-laplacian, def:pet-ch2-divergence-lie-derivative}
  \leanok
  \dcref{ch2:2.5.27}
  On an oriented Riemannian manifold:
  \begin{enumerate}
    \item compact support implies
      $\int_M\Delta f\,\mathrm{vol}=0$;
    \item
      $\operatorname{div}(fX)
        =g(\nabla f,X)+f\operatorname{div}X$;
    \item
      \[
        \Delta(f_1f_2)
          =(\Delta f_1)f_2+2g(\nabla f_1,\nabla f_2)
           +f_1\Delta f_2;
      \]
    \item for compactly supported $f_1$,
      $\int_Mf_1\Delta f_2\,\mathrm{vol}
       =-\int_Mg(\nabla f_1,\nabla f_2)\,\mathrm{vol}$;
    \item a compactly supported subharmonic or superharmonic function is
      constant; more generally, a subharmonic function attaining a global
      maximum, or a superharmonic function attaining a global minimum, is
      constant by the strong maximum principle.
  \end{enumerate}
  Parts (2)--(3) are formalized. Parts (1), (4), and (5) require Stokes'
  theorem and the maximum principle and are deferred.
  \sketch
  Apply the connection's Leibniz rule under the orthonormal trace; the two
  cross terms in the Laplacian combine to
  $2g(\nabla f_1,\nabla f_2)$.
\end{remark}

\begin{remark}[Incompressible vector fields]
  \label{rem:pet-ch2-ex-28}
  \lean{PetersenLib.exercise2_5_28}
  \source{petersen:page-0093}
  \uses{def:pet-ch2-divergence-lie-derivative, def:pet-ch2-covariant-derivative-tensor}
  \leanok
  \dcref{ch2:2.5.28}
  A vector field is \emph{incompressible} when
  $\operatorname{div}X=0$. Its local flows then preserve volume, and conversely.
  Every incompressible unit field on $\mathbb{R}^2$ is parallel, whereas
  $\mathbb{R}^3$ admits nonparallel examples. Namely,
  \[
    X(x)=\cos\langle e_3,x\rangle e_1
         +\sin\langle e_3,x\rangle e_2
  \]
  has unit length, $\nabla_{e_3}X|_0=e_2\ne0$, and zero covariant divergence.
  Indeed, with
  \[
    N(x)=-\sin\langle e_3,x\rangle e_1
          +\cos\langle e_3,x\rangle e_2,
  \]
  one has $\sum_i g(\nabla_{e_i}X,e_i)=\langle e_3,N\rangle=0$, which equals
  $\operatorname{div}X$ by
  \cref{prop:pet-ch2-divergence-trace-formula}.
  The example is formalized using
  $\sum_i g(\nabla_{e_i}X,e_i)=0$ rather than the volume-form definition, to
  avoid incompatible tangent-bundle inner-product instances in
  \cref{def:pet-ch2-divergence-lie-derivative}. The flow
  equivalence and the two-dimensional rigidity are deferred.
\end{remark}

\begin{remark}[Geodesic unit fields and integrability]
  \label{rem:pet-ch2-ex-29}
  \lean{PetersenLib.exercise2_5_29}
  \source{petersen:page-0093}
  \uses{def:pet-ch2-dual-one-form, def:pet-ch2-covariant-derivative-tensor}
  \leanok
  \dcref{ch2:2.5.29}
  Let $X$ be unit length with $\nabla_XX=0$.
  Then $X$ is locally a gradient exactly when $X^\perp$ is integrable. If
  $X^\perp$ has an integral submanifold through $p$, it is integrable near
  $p$; there are examples on $S^3$ for which $X$ is not a gradient.
  \sketch
  The formalized algebraic step is
  $L_X\theta_X=0$. Indeed,
  \[
    (L_X\theta_X)(Y)
      =D_Xg(X,Y)-g(X,[X,Y])
      =g(X,\nabla_YX)=0,
  \]
  using metric compatibility, torsion-freeness, $\nabla_XX=0$, and
  $D_Yg(X,X)=0$. Together with $\theta_X(X)=1$, this is the algebraic input
  for the Frobenius argument. The integrability conclusions and the $S^3$
  construction are deferred.
\end{remark}

\begin{remark}[Parallel complementary distributions]
  \label{rem:pet-ch2-ex-30}
  \lean{PetersenLib.exercise2_5_30}
  \source{petersen:page-0094}
  \uses{def:pet-ch2-covariant-derivative-tensor}
  \leanok
  \dcref{ch2:2.5.30}
  Let $E,F$ be orthogonal complementary parallel
  distributions. Each is integrable, and every point has a product
  neighborhood
  \[
    (U,g)=(V_E\times V_F,g|_{V_E}+g|_{V_F})
  \]
  by their integral submanifolds.
  \sketch
  The formalized algebraic part shows that any parallel distribution is
  involutive: torsion-freeness gives
  $[X,Y]=\nabla_XY-\nabla_YX$, whose two terms remain in the distribution.
  Frobenius integrability and the local metric splitting are deferred.
\end{remark}

\begin{remark}[Parallel vector fields and product metrics]
  \label{rem:pet-ch2-ex-31}
  \lean{PetersenLib.exercise2_5_31}
  \source{petersen:page-0094}
  \uses{def:pet-ch2-parallel-tensor, rem:pet-ch2-ex-30}
  \leanok
  \dcref{ch2:2.5.31}
  A parallel vector field $X$ has constant
  length. The line distribution generated by $X$ and its orthogonal complement
  are parallel, and locally
  \[
    (U,g)=(V\times I,g|_V+dt^2),
    \qquad TV\perp X.
  \]
  \sketch
  The formalized part is constancy of length:
  \[
    D_vg(X,X)
      =g(\nabla_vX,X)+g(X,\nabla_vX)=0.
  \]
  The parallel-distribution and product conclusions depend on
  \cref{rem:pet-ch2-ex-30} and are deferred.
\end{remark}

\begin{remark}[Leibniz rule for the inner product of tensors]
  \label{rem:pet-ch2-ex-32}
  \lean{PetersenLib.tensorFieldMetricInner_leibniz_orthonormal}
  \source{petersen:page-0094}
  \uses{def:pet-ch2-covariant-derivative-tensor, def:pet-ch2-tensor-algebra-inner-product}
  \leanok
  \dcref{ch2:2.5.32}
  For tensors $S,T$ of the same type,
  \[
    D_Xg(S,T)=g(\nabla_XS,T)+g(S,\nabla_XT).
  \]
  This is formalized for the tensor inner product computed in a smooth
  orthonormal frame at every point, using the frame conventions of
  \cref{rem:pet-ch2-ex-27} and \cref{rem:pet-ch2-ex-34}.
  \sketch
  The connection terms contributed by the
  frame cancel in pairs using
  $g(\nabla_XE_a,E_b)=-g(E_a,\nabla_XE_b)$. A normal-at-the-evaluation-point
  specialization, where these terms vanish individually, supplies the
  pointwise input to
  \cref{prop:pet-ch2-adjoint-integral-identity}.
\end{remark}

\begin{remark}[Almost complex structures, Nijenhuis tensor, K\"ahler condition]
  \label{rem:pet-ch2-ex-33}
  \lean{PetersenLib.exercise2_5_33}
  \source{petersen:page-0094}
  \uses{def:pet-ch2-lie-derivative-vector-field, def:pet-ch2-lie-derivative-tensor}
  \leanok
  \dcref{ch2:2.5.33}
  An almost complex structure is a $(1,1)$-tensor
  $J$ with $J^2=-I$. It is Hermitian when
  $g(JX,JY)=g(X,Y)$, and then
  $\omega(X,Y)=g(JX,Y)$ is its K\"ahler form. The Nijenhuis tensor is
  \[
    N(X,Y)=[JX,JY]-J[JX,Y]-J[X,JY]-[X,Y].
  \]
  It is tensorial; it vanishes for a complex structure, with the converse
  given by Newlander--Nirenberg. The form $\omega$ is alternating. Moreover
  $\nabla J=0$ implies $d\omega=0$, and the converse holds for a complex
  structure; such a Hermitian manifold is K\"ahler.
  Tensoriality of $N$ and alternation of $\omega$ are formalized; the
  integrability and exterior-derivative assertions are deferred.
  \sketch
  The bracket Leibniz terms cancel using $J^2=-I$, while
  $g(JX,JY)=g(X,Y)$ implies that $g(JX,Y)$ is alternating.
\end{remark}

\begin{remark}[Index notation for the covariant derivative]
  \label{rem:pet-ch2-ex-34}
  \lean{PetersenLib.exercise2_5_34}
  \source{petersen:page-0094}
  \uses{def:pet-ch2-covariant-derivative-index-notation, def:pet-ch2-divergence-lie-derivative, def:pet-ch2-covariant-derivative-adjoint}
  \leanok
  \dcref{ch2:2.5.34}
  With
  $\nabla_iT=\nabla_{\partial_i}T$ and
  $\nabla^iT=g^{ij}\nabla_jT$:
  \begin{enumerate}
    \item $df=\nabla_if\,dx^i$ and
      $\nabla f=\nabla^if\,\partial_i$;
    \item $(\nabla_iX)^i=\operatorname{div}X$;
    \item for a $(0,2)$-tensor,
      $(\nabla^iT)_{ij}=-(\nabla^*T)_j$.
  \end{enumerate}
  All three identities are formalized in a positively oriented orthonormal
  frame $E_i$ with dual coframe $\sigma^i=g(E_i,\cdot)$.
  \sketch
  The first is the
  orthonormal expansion of $\nabla f$, the second is
  \cref{prop:pet-ch2-divergence-trace-formula}, and the third is the definition
  of \cref{def:pet-ch2-covariant-derivative-adjoint}.
\end{remark}
